\chapter{Strategic Mealy Systems and Open Game Semantics}
\label{chap:strategic-mealy-games}

\section{Orientation}
\label{sec:strategic-mealy-orientation}

The preceding chapters developed a categorical systems theory in which internal categories are
interfaces, Mealy morphisms are stateful open systems, Mealy simulations compare implementations,
downstream actions are environments, residual semantics extracts behaviour, product tensor describes
parallel execution, and coproduct trace closes hidden event-flow interfaces. The purpose of the
present chapter is to construct strategic interaction inside that theory without discarding its
\(2\)-dimensional structure.

Throughout the chapter, the ambient category is a Grothendieck topos. This permits the strategic
constructions to be carried out internally. Strategy spaces are objects rather than external sets,
best-response relations are subobjects, truth values lie in the subobject classifier, finite
counterfactual query words are internal lists, and the internally defined constructions are stable
under arbitrary base change up to canonical coherent isomorphism. The exactness and local cartesian
closure required by the residual semantics of Chapter~3 are automatic. The stable finite and
countable coproducts required by the finite-path trace of Chapter~9 are also available.

The chapter has three parts.

The first part constructs strategic Mealy game theory. It begins by strengthening initialization so
that selected states are closed under input processing. This strengthening is forced by horizontal
composition of output-lax Mealy simulations: a downstream selected state must be transported along
the output-comparison arrow of the upstream simulation. Operational bidirectional arenas are then
constructed as Mealy optics. Residual reparametrizers carry coherent discard sections, and pure
map-representable residual slides are identified by a homwise coidentifier. Frozen strategy
parameters are represented by discrete internal categories. Strategic contexts are supplied by an
internal context doctrine, and counterfactual response is represented by an initialized internal
Mealy machine over an internal query category. These data form a symmetric monoidal bicategory of
response-strict strategic Mealy games.

The second part proves two bridge theorems. The first identifies internal deterministic open games
with open-normal strategic Mealy games. The second identifies open-normal strategic Mealy games
with structured canonical branch-faithful protocol machines. Thus ordinary play, coplay, and
best-response semantics are recovered as a rigid normal form of the general strategic Mealy theory,
while the canonical protocol compiler places the forward, backward, and response branches inside one
initialized typed machine.

The third part gives residual, execution, and feedback semantics. There are two related execution
constructions.

First, primitive relative execution of a genuine Mealy pipeline factors through a visible
intermediate handshake object. The composite primitive span is recovered by tracing over that
handshake object. Persistent implementation state remains in the visible relative execution
configuration.

Second, a stateless object map between discrete interfaces is elaborated into a genuine event of a
free internal step category. Restricting Chapter~6 primitive execution to the length-one event
generators recovers the graph span of the map. This step elaboration turns the forward and backward
maps of an open-normal game into genuine primitive execution spans. Sequential arena composition is
then recovered as the finite-path trace of a private linear subdivision.

The bridge theorems and the feedback theorem may be displayed as
\[
  \boxed{
  \mathbf{OpenGame}_2(\mathcal E)
  \cong
  \mathbf{OGMly}^{\mathrm{nf}}_2(\mathcal E)
  \cong
  \mathbf{CanProt}^{\mathrm{str}}_2(\mathcal E)
  }
\]
and
\[
  \boxed{
  \operatorname{Exec}_{\mathrm{step}}
  \bigl(
    \mathsf{Prot}_{\mathrm{ar}}
    (\mathcal H\circ\mathcal G)
  \bigr)
  \cong
  \operatorname{Tr}^{H}
  \bigl(
    \widetilde U_{\mathcal H,\mathcal G}
  \bigr).
  }
\]

The stateful pipeline factorization is
\[
  \boxed{
  m_{Q\circ P,Y}
  \cong
  \operatorname{Tr}^{H_{Q,P;Y}}
  \bigl(
    \widetilde m_{Q,P;Y}
  \bigr).
  }
\]

\subsection{The strategic systems dictionary}
\label{subsec:strategic-systems-dictionary}

\begin{center}
\begin{tabular}{p{0.29\textwidth}|p{0.61\textwidth}}
\textbf{Open-game notion}
&
\textbf{Strategic Mealy interpretation}
\\
\hline
\((X,S)\to(Y,R)\)
&
A signed boundary with positive source \(X\), negative source \(S\), positive target \(Y\), and
negative target \(R\).
\\[1mm]
\(\Sigma\)
&
A frozen internal strategy parameter indexing controller implementations.
\\[1mm]
\(P:\Sigma\times X\to Y\)
&
The visible positive output of the forward phase.
\\[1mm]
\(C:\Sigma\times X\times R\to S\)
&
The visible negative output of the backward phase, using data retained from the forward phase.
\\[1mm]
\(k:Y\to R\)
&
A stationary continuation, represented internally by an element of the exponential \(R^Y\).
\\[1mm]
\((x,k)\)
&
An upstream point together with a stationary downstream continuation.
\\[1mm]
Counterfactual deviation
&
A typed query asking the response machine to evaluate an alternative strategy parameter.
\\[1mm]
Best response
&
A subobject of
\(\Sigma\times\operatorname{Ctx}\times\Sigma\), equivalently a classifier into \(\Omega\).
\\[1mm]
Selection function
&
A response controller consuming a counterfactual outcome profile and classifying acceptable
alternatives.
\\[1mm]
Equilibrium
&
A diagonal response: the incumbent is accepted as a response to itself.
\\[1mm]
Sequential composition
&
Residual optic interconnection together with coherent decomposition of the outer context into local
contexts.
\\[1mm]
Parallel composition
&
Product tensor together with partial evaluation of the joint continuation against the other
incumbent component.
\\[1mm]
Strategy reparametrization
&
A \(2\)-cell combining a contravariant map of strategy objects, an optic simulation, and a
response-strict evaluator simulation.
\\[1mm]
Protocol implementation
&
One initialized Mealy morphism with stable forward, backward, and evaluator carrier branches and an
explicit strategy port.
\\[1mm]
Relative execution
&
Primitive execution spans on visible implementation configurations.
\\[1mm]
Step elaboration
&
A free internal event category whose primitive generator spans recover graph spans of stateless
maps.
\\[1mm]
Feedback
&
Finite-path closure of private execution routes after relative execution evaluation.
\end{tabular}
\end{center}

Three distinctions are maintained throughout.

First, an optic residual interface is not the implementation-state carrier of either Mealy leg. The
residual is exposed by the forward leg and consumed by the backward leg. Each leg may additionally
carry persistent implementation state.

Second, product-monoidal residual composition and feedback perform different operations. Product
tensor records simultaneously available data. Feedback hides event-flow routes only after a phase
protocol has made causal order explicit.

Third, the protocol input and output categories of a compiled game are implementation artefacts.
They are not the external \(0\)-cells of the canonical protocol bicategory. Canonical protocol
composition is compiler composition obtained from extracted strategic data, not ordinary Mealy
pipeline composition of the underlying compiled machines.

\section{Standing framework and notation}
\label{sec:strategic-standing-framework}

Throughout this chapter, let \(\mathcal E\) be a Grothendieck topos with chosen pullbacks. All
objects, maps, relations, categories, functors, natural transformations, Mealy morphisms, and Mealy
simulations are internal to \(\mathcal E\), unless explicitly stated otherwise.

A Grothendieck topos is Barr-exact and locally cartesian closed. Hence it satisfies
Assumption~3.1.1. It also has disjoint pullback-stable small coproducts. In particular, it is a
feedback span base in the sense of Chapter~9 for every finite and countable coproduct used below.

Generalized-element notation abbreviates morphism equations in \(\mathcal E\). Consequently a
variable
\[
  \sigma:\Sigma
\]
denotes an internal variable, or equivalently a generalized element after arbitrary base change. No
well-pointedness condition is used.

Fix an external universe in which the internal categories and the relevant hom-categories are small.
Fix also:

\begin{enumerate}
  \item chosen finite and countable coproducts;
  \item a chosen free-internal-category functor on internal graphs;
  \item chosen pushouts used to adjoin protocol generators;
  \item chosen effective quotients of internal congruences;
  \item chosen homwise coidentifiers and connected-component quotients in the external universe.
\end{enumerate}

These choices allow each construction in a fixed base to be represented by a specified object rather
than merely by an isomorphism class.

For every map
\[
  g:V\to U,
\]
the pullback functor
\[
  g^*:\mathcal E/U\longrightarrow\mathcal E/V
\]
has both adjoints
\[
  \Sigma_g\dashv g^*\dashv\Pi_g.
\]
It therefore preserves finite limits and small colimits. It also preserves regular epimorphisms,
monomorphisms, regular images, kernel pairs, and effective quotients. Beck--Chevalley gives the
canonical comparison isomorphisms required for slice exponentials and dependent products.

Consequently there are canonical coherent isomorphisms
\[
  g^*(X^\ast)\cong(g^*X)^\ast,
\]
\[
  g^*\mathsf{FreeCat}(G)
  \cong
  \mathsf{FreeCat}(g^*G),
\]
\[
  g^*(\mathcal C^\simeq)
  \cong
  (g^*\mathcal C)^\simeq,
\]
\[
  g^*\operatorname{Reach}(P,i)
  \cong
  \operatorname{Reach}(g^*P,g^*i),
\]
\[
  g^*\operatorname{Der}_0(k)
  \cong
  \operatorname{Der}_0(g^*k),
\]
\[
  g^*(D_k/{\equiv_k})
  \cong
  D_{g^*k}/{\equiv_{g^*k}},
\]
and
\[
  g^*\operatorname{Tr}^{H}(R)
  \cong
  \operatorname{Tr}^{g^*H}(g^*R).
\]
These comparison isomorphisms are pseudonatural in \(g\). The external connected-component and
homwise coidentifier quotients are not included in this base-change assertion.

Write
\[
  \mathcal B:=\operatorname{Mly}(\mathcal E).
\]
Its objects are internal categories, its \(1\)-cells are Mealy morphisms, and its \(2\)-cells are
Mealy simulations. Finite product induces a symmetric monoidal bicategory
\[
  (\mathcal B,\times,\mathbbm1).
\]

For a Mealy morphism
\[
  P:A\rightsquigarrow B,
\]
write
\[
  p_P:P_0\to A_0,
  \qquad
  q_P:P_0\to B_0
\]
for its carrier anchors,
\[
  \delta_P:
  A_1\times_{t_A,A_0,p_P}P_0\to P_0
\]
for its transition map, and
\[
  \omega_P:
  A_1\times_{t_A,A_0,p_P}P_0\to B_1
\]
for its output map.

The target-to-source convention gives
\[
  p_P\delta_P(a,x)=s_A(a),
\]
\[
  t_B\omega_P(a,x)=q_P(x),
  \qquad
  s_B\omega_P(a,x)=q_P\delta_P(a,x).
\]
For composable
\[
  a:u\to v,
  \qquad
  b:v\to w,
  \qquad
  p_P(x)=w,
\]
the multiplication equations are
\[
  \delta_P(b\circ a,x)
  =
  \delta_P(a,\delta_P(b,x))
\]
and
\[
  \omega_P(b\circ a,x)
  =
  \omega_P(b,x)
  \circ
  \omega_P(a,\delta_P(b,x)).
\]

The displayed direction
\[
  \Theta:P\Longrightarrow Q
\]
of a Mealy simulation is opposite to the direction of its carrier comparison
\[
  \theta:Q_0\longrightarrow P_0.
\]
Write
\[
  \alpha_x:
  q_P(\theta x)\longrightarrow q_Q(x)
\]
for its output-comparison component. Its state and output equations are
\[
  \theta\delta_Q(a,x)
  =
  \delta_P(a,\theta x)
\]
and
\[
  \alpha_x\circ\omega_P(a,\theta x)
  =
  \omega_Q(a,x)
  \circ
  \alpha_{\delta_Q(a,x)}.
\]

For an internal functor
\[
  F:A\to B,
\]
write
\[
  F_\#:A\rightsquigarrow B
\]
for its map-representable Mealy morphism.

We write \(X+Y\) for coproducts of objects and for the coproduct monoidal tensor of spans. We reserve
\[
  R\boxplus S
\]
for the homwise coproduct of parallel spans with common source and target.

We suppress specified associators, unitors, tensorators, and canonical pullback-comparison
isomorphisms when their placement is uniquely forced by typing. No non-coherence \(2\)-cell is
suppressed.

\section{Transition-closed initialized Mealy morphisms}
\label{sec:initialized-mealy-morphisms}

An arbitrary family of selected states is not stable under horizontal composition of output-lax
Mealy simulations. The selected state of the downstream machine must be transported along the
output-comparison arrow of the upstream simulation. Initialization must therefore be closed under
input processing.

\begin{definition}[Transition-closed initialization]
\label{def:transition-closed-initialization}
Let
\[
  P:A\rightsquigarrow B
\]
be a Mealy morphism. A transition-closed initialization of \(P\) consists of:
\begin{enumerate}
  \item an object \(J_P\) and an anchor
  \[
    r_P:J_P\longrightarrow A_0;
  \]
  \item a right internal \(A\)-action
  \[
    \lambda_P:
    A_1\times_{t_A,A_0,r_P}J_P
    \longrightarrow
    J_P;
  \]
  \item an equivariant map
  \[
    j_P:J_P\longrightarrow P_0
  \]
  satisfying
  \[
    p_Pj_P=r_P
  \]
  and
  \[
    j_P\lambda_P(a,i)
    =
    \delta_P(a,j_Pi).
  \]
\end{enumerate}
The action laws are
\[
  \lambda_P(1_{r_P(i)},i)=i
\]
and
\[
  \lambda_P(a,\lambda_P(b,i))
  =
  \lambda_P(b\circ a,i)
\]
whenever \(a:u\to v\), \(b:v\to w\), and \(r_P(i)=w\).

The resulting initialized Mealy morphism is written
\[
  \mathbf P=(P,J_P,\lambda_P,j_P).
\]
\end{definition}

A ready family for \(\mathbf P\) is additional data
\[
  \rho_P:R_P\longrightarrow J_P.
\]
It need not be transition-closed. Transition-closed initialization supports composition; the ready
family specifies the states whose observations are to be exposed, generated, or minimized.

\begin{definition}[Total initialization]
\label{def:total-initialization}
The total initialization of \(P\) is
\[
  \operatorname{Tot}(P)
  :=
  (P,P_0,\delta_P,1_{P_0}),
\]
with initialization anchor \(p_P:P_0\to A_0\).
\end{definition}

\begin{definition}[Initialized Mealy simulation]
\label{def:initialized-mealy-simulation}
Let
\[
  \mathbf P=(P,J_P,\lambda_P,j_P),
  \qquad
  \mathbf Q=(Q,J_Q,\lambda_Q,j_Q)
\]
be transition-closed initialized Mealy morphisms \(A\rightsquigarrow B\).

An initialized Mealy simulation
\[
  (u,\Theta):\mathbf P\Longrightarrow\mathbf Q
\]
consists of a Mealy simulation
\[
  \Theta:P\Longrightarrow Q
\]
with carrier comparison
\[
  \theta:Q_0\to P_0,
\]
together with a morphism of right \(A\)-actions
\[
  u:J_Q\longrightarrow J_P
\]
such that
\[
  j_Pu=\theta j_Q.
\]
Thus
\[
  r_Pu=r_Q
\]
and
\[
  u\lambda_Q(a,k)
  =
  \lambda_P(a,uk).
\]
\end{definition}

\begin{definition}[Strict initialized simulation]
\label{def:strict-initialized-simulation}
An initialized Mealy simulation
\[
  (u,\Theta):\mathbf P\Longrightarrow\mathbf Q
\]
is strict when the output-comparison component of \(\Theta\) is everywhere an identity arrow.
\end{definition}

\subsection{Pipeline initialization}

Let
\[
  \mathbf P:A\rightsquigarrow B,
  \qquad
  \mathbf Q:B\rightsquigarrow C.
\]
Define
\[
  J_{Q\circ P}
  :=
  J_P\times_{q_Pj_P,B_0,r_Q}J_Q.
\]
Its carrier map is
\[
  j_{Q\circ P}(i,k)
  :=
  (j_Pi,j_Qk),
\]
and its input action is
\[
  \lambda_{Q\circ P}(a,(i,k))
  :=
  \Bigl(
    \lambda_P(a,i),
    \lambda_Q(\omega_P(a,j_Pi),k)
  \Bigr).
\]

\begin{lemma}[Pipeline initialization]
\label{lem:pipeline-initialization}
The preceding formulas define a transition-closed initialization of \(Q\circ P\).
\end{lemma}

\begin{proof}
Since \((i,k)\) lies in the pullback,
\[
  q_Pj_Pi=r_Qk.
\]
The arrow
\[
  \omega_P(a,j_Pi)
\]
has target \(q_Pj_Pi\), so it acts on \(k\). Its source is
\[
\begin{aligned}
  s_B\omega_P(a,j_Pi)
  &=
  q_P\delta_P(a,j_Pi)
  \\
  &=
  q_Pj_P\lambda_P(a,i).
\end{aligned}
\]
Hence the displayed pair again lies in the pullback.

Equivariance is
\[
\begin{aligned}
  j_{Q\circ P}
  \lambda_{Q\circ P}(a,(i,k))
  &=
  \Bigl(
    \delta_P(a,j_Pi),
    \delta_Q(\omega_P(a,j_Pi),j_Qk)
  \Bigr)
  \\
  &=
  \delta_{Q\circ P}
  (a,(j_Pi,j_Qk)).
\end{aligned}
\]

The identity law follows from the identity laws of \(P\), \(Q\), and the two actions. For composable
\(a:u\to v\) and \(b:v\to w\), the first component is the action multiplication law for
\(\lambda_P\). For the second component, the Mealy multiplication law gives
\[
  \omega_P(b,j_Pi)
  \circ
  \omega_P(a,\delta_P(b,j_Pi))
  =
  \omega_P(b\circ a,j_Pi).
\]
The action multiplication law for \(\lambda_Q\) then gives the result.
\end{proof}

\subsection{Horizontal composition of initialized simulations}

Let
\[
  (u,\Theta):\mathbf P\Longrightarrow\mathbf P',
  \qquad
  (v,\Xi):\mathbf Q\Longrightarrow\mathbf Q'
\]
with
\[
  P,P':A\rightsquigarrow B,
  \qquad
  Q,Q':B\rightsquigarrow C.
\]
Write
\[
  \theta:P'_0\to P_0,
  \qquad
  \kappa:Q'_0\to Q_0
\]
for the carrier comparisons, and write
\[
  \alpha_x:
  q_P(\theta x)\longrightarrow q_{P'}(x)
\]
for the output comparison of \(\Theta\).

Define
\[
\begin{aligned}
  h_{\Xi,\Theta}:
  J_{Q'\circ P'}
  &\longrightarrow
  J_{Q\circ P},
  \\
  (i',k')
  &\longmapsto
  \Bigl(
    ui',
    \lambda_Q(\alpha_{j_{P'}i'},vk')
  \Bigr).
\end{aligned}
\]

\begin{lemma}[Horizontal initialization comparison]
\label{lem:horizontal-initialization-comparison}
The map \(h_{\Xi,\Theta}\) is a morphism of right \(A\)-actions and factors the carrier comparison of
the horizontal Mealy simulation
\[
  \Xi\ast\Theta:
  Q\circ P
  \Longrightarrow
  Q'\circ P'.
\]
\end{lemma}

\begin{proof}
Since
\[
  t_B\alpha_{j_{P'}i'}
  =
  q_{P'}j_{P'}i'
  =
  r_{Q'}k'
  =
  r_Q(vk'),
\]
the second component is defined. Its resulting anchor is
\[
\begin{aligned}
  r_Q\lambda_Q(\alpha_{j_{P'}i'},vk')
  &=
  s_B\alpha_{j_{P'}i'}
  \\
  &=
  q_P\theta j_{P'}i'
  \\
  &=
  q_Pj_Pui'.
\end{aligned}
\]
Thus the pair lies in \(J_{Q\circ P}\).

Using
\[
  j_Pu=\theta j_{P'},
  \qquad
  j_Qv=\kappa j_{Q'},
\]
one obtains
\[
  j_{Q\circ P}h_{\Xi,\Theta}(i',k')
  =
  \Bigl(
    \theta j_{P'}i',
    \delta_Q
    \bigl(
      \alpha_{j_{P'}i'},
      \kappa j_{Q'}k'
    \bigr)
  \Bigr),
\]
which is the carrier comparison of \(\Xi\ast\Theta\).

For equivariance, let \(a\) be an input arrow for \(P'\). Processing first by
\(h_{\Xi,\Theta}\) and then by \(a\) gives the second component
\[
  \lambda_Q
  \bigl(
    \alpha_{j_{P'}i'}
    \circ
    \omega_P(a,\theta j_{P'}i'),
    vk'
  \bigr).
\]
Processing first by \(a\) and then by \(h_{\Xi,\Theta}\) gives
\[
  \lambda_Q
  \bigl(
    \omega_{P'}(a,j_{P'}i')
    \circ
    \alpha_{\delta_{P'}(a,j_{P'}i')},
    vk'
  \bigr).
\]
These are equal by the output-coherence equation for \(\Theta\). The first components agree by
equivariance of \(u\).
\end{proof}

The identity initialized Mealy morphism on \(A\) is the identity Mealy morphism with total
initialization. Tensor is defined coordinatewise:
\[
  \mathbf P\times\mathbf Q
  :=
  (P\times Q,J_P\times J_Q,\lambda_P\times\lambda_Q,j_P\times j_Q).
\]

\begin{theorem}[Initialized Mealy bicategory]
\label{thm:initialized-mealy-bicategory}
Transition-closed initialized Mealy morphisms and initialized Mealy simulations form a bicategory
\[
  \operatorname{IMly}^{\mathrm{cl}}(\mathcal E).
\]
Finite product induces a symmetric monoidal bicategory structure. Forgetting initialization defines
a symmetric monoidal pseudofunctor
\[
  U:
  \operatorname{IMly}^{\mathrm{cl}}(\mathcal E)
  \longrightarrow
  \operatorname{Mly}(\mathcal E).
\]
Total initialization defines a symmetric monoidal pseudofunctor
\[
  \operatorname{Tot}:
  \operatorname{Mly}(\mathcal E)
  \longrightarrow
  \operatorname{IMly}^{\mathrm{cl}}(\mathcal E)
\]
with
\[
  U\operatorname{Tot}=1.
\]
\end{theorem}

\begin{proof}
Vertical composition composes the action maps. If
\[
  (u,\Theta):\mathbf P\Longrightarrow\mathbf Q,
  \qquad
  (v,\Xi):\mathbf Q\Longrightarrow\mathbf R,
\]
then
\[
  j_Puv
  =
  \theta j_Qv
  =
  \theta\kappa j_R,
\]
so \(uv\) factors the composite carrier comparison.

Horizontal composition is supplied by
Lemma~\ref{lem:horizontal-initialization-comparison}. Expanding either parenthesization of a triple
horizontal composite produces the same ordered sequence of output-comparison arrows acting on the
final initialization state. Equality follows from the action multiplication law and the
associativity equation for horizontal Mealy simulation.

The associators and unitors are the canonical equivariant isomorphisms between the corresponding
iterated pullbacks. Their pentagon and triangle equations follow from the equations in
\(\operatorname{Mly}(\mathcal E)\) and uniqueness of the canonical pullback comparisons.

Product acts coordinatewise. Its coherence is inherited from finite-product coherence in
\(\mathcal E\).

For a Mealy simulation
\[
  \Theta:P\Longrightarrow Q
\]
with carrier comparison \(\theta:Q_0\to P_0\), total initialization uses \(\theta\) as its
initialization map. The transition equation for \(\Theta\) is exactly the required equivariance
equation. The canonical pullback comparison identifies total initialization of a pipeline with the
pipeline of the total initializations. The identity and tensor comparisons are immediate.
\end{proof}

\section{Discardable residual reparametrizers}
\label{sec:discardable-residual-reparametrizers}

An arbitrary stateful residual reparametrizer cannot be discarded functorially. For
\[
  R:N\rightsquigarrow M,
\]
the composite
\[
  (!_M)_\#\circ R:N\rightsquigarrow\mathbbm1
\]
has carrier \(R_0\), whereas
\[
  (!_N)_\#:N\rightsquigarrow\mathbbm1
\]
has carrier \(N_0\). A simulation in the displayed direction requires a coherent \(R\)-state over
every object of \(N\).

\begin{definition}[Discardable residual reparametrizer]
\label{def:discardable-residual-reparametrizer}
A discardable residual reparametrizer
\[
  (R,e_R):N\rightsquigarrow M
\]
consists of a Mealy morphism \(R:N\rightsquigarrow M\) and a map
\[
  e_R:N_0\longrightarrow R_0
\]
such that
\[
  p_Re_R=1_{N_0}
\]
and
\[
  \delta_R(a,e_R(v))
  =
  e_R(u)
\]
for every arrow \(a:u\to v\) of \(N\).

Equivalently, \(e_R\) is the carrier map of a strict initialized Mealy simulation
\[
  \epsilon_R:
  \operatorname{Tot}((!_M)_\#\circ R)
  \Longrightarrow
  \operatorname{Tot}((!_N)_\#).
\]
\end{definition}

\begin{definition}[Discard-preserving simulation]
\label{def:discard-preserving-simulation}
A Mealy simulation
\[
  \rho:(R,e_R)\Longrightarrow(R',e_{R'})
\]
is discard-preserving when its carrier comparison
\[
  \rho_0:R'_0\to R_0
\]
satisfies
\[
  \rho_0e_{R'}=e_R.
\]
\end{definition}

\begin{proposition}[Closure of discardable reparametrizers]
\label{prop:discardable-closure}
Discardable residual reparametrizers are closed under identities, pipeline composition, and product
tensor. Discard-preserving simulations are closed under vertical composition, horizontal
composition, and tensor.
\end{proposition}

\begin{proof}
The identity reparametrizer has
\[
  e_{1_N}=1_{N_0}.
\]

For
\[
  S:L\rightsquigarrow N,
  \qquad
  R:N\rightsquigarrow M,
\]
define
\[
  e_{R\circ S}(\ell)
  :=
  \bigl(
    e_S(\ell),
    e_R(q_Se_S(\ell))
  \bigr)
\]
in the pullback carrier
\[
  S_0\times_{N_0}R_0.
\]
Let \(a:u\to v\) be an arrow of \(L\), and put
\[
  \beta:=\omega_S(a,e_S(v)).
\]
Discardability of \(S\) gives
\[
  \delta_S(a,e_S(v))=e_S(u).
\]
The arrow \(\beta\) has target \(q_Se_S(v)\) and source \(q_Se_S(u)\). Hence
\[
  \delta_R
  \bigl(
    \beta,
    e_R(q_Se_S(v))
  \bigr)
  =
  e_R(q_Se_S(u)).
\]
Therefore
\[
  \delta_{R\circ S}
  (a,e_{R\circ S}(v))
  =
  e_{R\circ S}(u).
\]

For tensor, define
\[
  e_{R\times S}(n,\ell)
  :=
  (e_R(n),e_S(\ell)).
\]

Vertical composition of discard-preserving simulations is immediate. For horizontal composition,
let
\[
  \rho:S\Longrightarrow S',
  \qquad
  \sigma:R\Longrightarrow R'
\]
be discard-preserving, and write
\[
  \alpha_x:
  q_S(\rho_0x)\longrightarrow q_{S'}(x)
\]
for the output comparison of \(\rho\). The horizontal carrier comparison sends
\[
  e_{R'\circ S'}(\ell)
  =
  \bigl(
    e_{S'}(\ell),
    e_{R'}(q_{S'}e_{S'}(\ell))
  \bigr)
\]
to
\[
  \left(
    e_S(\ell),
    \delta_R
    \bigl(
      \alpha_{e_{S'}(\ell)},
      e_R(q_{S'}e_{S'}(\ell))
    \bigr)
  \right).
\]
The displayed transition is defined because
\[
  t_N\alpha_{e_{S'}(\ell)}
  =
  q_{S'}e_{S'}(\ell).
\]
Its source is
\[
  s_N\alpha_{e_{S'}(\ell)}
  =
  q_Se_S(\ell).
\]
Discardability of \(R\) therefore gives
\[
  \delta_R
  \bigl(
    \alpha_{e_{S'}(\ell)},
    e_R(q_{S'}e_{S'}(\ell))
  \bigr)
  =
  e_R(q_Se_S(\ell)).
\]
Thus the resulting state is
\[
  \bigl(
    e_S(\ell),
    e_R(q_Se_S(\ell))
  \bigr)
  =
  e_{R\circ S}(\ell).
\]
Tensor preservation is coordinatewise.
\end{proof}

Every map-representable residual change
\[
  h_\#:N\rightsquigarrow M
\]
is discardable. Its carrier is \(N_0\), and one takes
\[
  e_{h_\#}=1_{N_0}.
\]

\section{Operational Mealy optics}
\label{sec:operational-mealy-optics}

\begin{definition}[Signed interface]
\label{def:signed-interface}
A signed interface is a pair
\[
  A=(A^+,A^-)
\]
of internal categories.
\end{definition}

\begin{definition}[Mealy optic presentation]
\label{def:mealy-optic-presentation}
For signed interfaces
\[
  A=(A^+,A^-),
  \qquad
  B=(B^+,B^-),
\]
a Mealy optic presentation
\[
  \mathcal A:A\rightsquigarrow_{\mathrm{opt}}B
\]
consists of an internal category \(M\), called its residual interface, and transition-closed
initialized Mealy morphisms
\[
  \mathbf F:A^+\rightsquigarrow M\times B^+,
\]
\[
  \mathbf G:M\times B^-\rightsquigarrow A^-.
\]
We write
\[
  \mathcal A=(M,\mathbf F,\mathbf G).
\]
\end{definition}

\subsection{Operational optic simulations}

Let
\[
  \mathcal A=(M,\mathbf F,\mathbf G),
  \qquad
  \mathcal A'=(N,\mathbf F',\mathbf G')
\]
have the same signed source and target.

\begin{definition}[Raw operational optic simulation]
\label{def:raw-operational-optic-simulation}
A raw operational optic simulation
\[
  (R,e_R,\alpha,\beta):
  \mathcal A\Longrightarrow\mathcal A'
\]
consists of:
\begin{enumerate}
  \item a discardable residual reparametrizer
  \[
    (R,e_R):N\rightsquigarrow M;
  \]
  \item an initialized Mealy simulation
  \[
    \alpha:
    \mathbf F
    \Longrightarrow
    (\operatorname{Tot}(R)\times1_{B^+})\circ\mathbf F';
  \]
  \item an initialized Mealy simulation
  \[
    \beta:
    \mathbf G\circ(\operatorname{Tot}(R)\times1_{B^-})
    \Longrightarrow
    \mathbf G'.
  \]
\end{enumerate}
\end{definition}

\begin{definition}[Optic modification]
\label{def:optic-modification}
Let
\[
  (R,e_R,\alpha,\beta),
  \qquad
  (R',e_{R'},\alpha',\beta')
\]
be raw operational optic simulations
\(\mathcal A\Longrightarrow\mathcal A'\).

A modification
\[
  \rho:
  (R,e_R,\alpha,\beta)
  \Rrightarrow
  (R',e_{R'},\alpha',\beta')
\]
is a discard-preserving Mealy simulation
\[
  \rho:R\Longrightarrow R'
\]
such that
\[
  \alpha'
  =
  \bigl(
    (\operatorname{Tot}(\rho)\times1_{B^+})\ast\mathbf F'
  \bigr)
  \circ\alpha
\]
and
\[
  \beta
  =
  \beta'
  \circ
  \bigl(
    \mathbf G\ast(\operatorname{Tot}(\rho)\times1_{B^-})
  \bigr).
\]
\end{definition}

For fixed \(\mathcal A,\mathcal A'\), raw simulations and modifications form a category. An
operational optic simulation
\[
  \mathcal A\Longrightarrow\mathcal A'
\]
is a connected component of this category.

Thus residual reparametrizers and their represented leg data are retained only modulo the
modification equations and the resulting zigzags. Stateful residual reparametrization is not
discarded merely because pure presentation changes will later be coidentified.

\subsection{Composition}

If
\[
  (R,e_R,\alpha,\beta):\mathcal A\Longrightarrow\mathcal A',
\]
\[
  (S,e_S,\gamma,\delta):\mathcal A'\Longrightarrow\mathcal A'',
\]
their vertical composite uses the residual pipeline \(R\circ S\). Its forward cell is the pasting
\[
  \mathbf F
  \overset{\alpha}{\Longrightarrow}
  (R\times1)\mathbf F'
  \overset{R\ast\gamma}{\Longrightarrow}
  (R\times1)(S\times1)\mathbf F''
  \Longrightarrow
  ((R\circ S)\times1)\mathbf F'',
\]
and its backward cell is the dual pasting.

For presentations
\[
  \mathcal A=(M,\mathbf F,\mathbf G):A\rightsquigarrow_{\mathrm{opt}}B,
\]
\[
  \mathcal B=(N,\mathbf H,\mathbf K):B\rightsquigarrow_{\mathrm{opt}}C,
\]
define
\[
  \mathcal B\circ_{\mathrm{opt}}\mathcal A
  :=
  \left(
    M\times N,
    (1_M\times\mathbf H)\circ\mathbf F,
    \mathbf G\circ(1_M\times\mathbf K)
  \right).
\]
The identity optic has residual \(\mathbbm1\) and the product unitors.

Horizontal composition of optic simulations uses the residual reparametrizer \(R\times S\), together
with the monoidal-pseudofunctor interchangers of product tensor on initialized Mealy morphisms.

\begin{theorem}[Operational Mealy optic bicategory]
\label{thm:operational-mealy-optic-bicategory}
Signed internal interfaces, Mealy optic presentations, and operational optic simulations form a
bicategory
\[
  \operatorname{MOpt}_{\mathrm{op}}(\mathcal E).
\]
Product tensor induces a symmetric monoidal bicategory structure
\[
  \bigl(
    \operatorname{MOpt}_{\mathrm{op}}(\mathcal E),
    \otimes,
    (\mathbbm1,\mathbbm1)
  \bigr).
\]
\end{theorem}

\begin{proof}
Raw vertical composition is associative up to the associator of
\(\operatorname{IMly}^{\mathrm{cl}}(\mathcal E)\). The residual associator is a
discard-preserving map-representable simulation, and the two leg pastings satisfy the modification
equations. Vertical composition is therefore associative after passage to connected components.

Horizontal composition is well defined because product is a pseudofunctor on initialized Mealy
morphisms and discardable residual reparametrizers are closed under product. The tensor of two
modifications, followed by the forced interchanger pastings, is a modification of the corresponding
horizontal composites.

Interchange follows from interchange in
\(\operatorname{IMly}^{\mathrm{cl}}(\mathcal E)\). The associator, unitors, and symmetry optics use
map-representable residual changes induced by finite-product coherence maps. Their coherence
equations reduce to those in the symmetric monoidal bicategory
\(\operatorname{IMly}^{\mathrm{cl}}(\mathcal E)\).
\end{proof}

\subsection{Visible output and continuation closure}

Let
\[
  \mathcal A=(M,\mathbf F,\mathbf G):A\rightsquigarrow_{\mathrm{opt}}B.
\]
Define its visible positive process by
\[
  \operatorname{out}(\mathcal A)
  :=
  \bigl(
    (!_M)_\#\times1_{B^+}
  \bigr)
  \circ\mathbf F.
\]
For an initialized continuation
\[
  \mathbf K:B^+\rightsquigarrow B^-,
\]
define
\[
  \operatorname{cl}_{\mathbf K}(\mathcal A)
  :=
  \mathbf G\circ(1_M\times\mathbf K)\circ\mathbf F.
\]

\begin{lemma}[Functoriality in the continuation]
\label{lem:continuation-functoriality}
For every operational optic \(\mathcal A=(M,\mathbf F,\mathbf G)\), continuation closure defines a
functor
\[
  \operatorname{Cl}_{\mathcal A}:
  \operatorname{IMly}^{\mathrm{cl}}(\mathcal E)(B^+,B^-)
  \longrightarrow
  \operatorname{IMly}^{\mathrm{cl}}(\mathcal E)(A^+,A^-).
\]
On a continuation simulation
\[
  \zeta:\mathbf K\Longrightarrow\mathbf K',
\]
its value is
\[
  \mathbf G\ast(1_M\times\zeta)\ast\mathbf F.
\]
Likewise, postcomposition with \(\operatorname{out}(\mathcal A)\) defines a functor
\[
  \operatorname{IMly}^{\mathrm{cl}}(\mathcal E)(\mathbbm1,A^+)
  \longrightarrow
  \operatorname{IMly}^{\mathrm{cl}}(\mathcal E)(\mathbbm1,B^+).
\]
\end{lemma}

\begin{proof}
Both assignments are whiskering operations in the bicategory
\(\operatorname{IMly}^{\mathrm{cl}}(\mathcal E)\). Their identity and composition laws are the
bicategorical unitor and associator equations.
\end{proof}

\begin{proposition}[Output and closure comparisons]
\label{prop:output-closure-comparisons}
A raw operational optic simulation
\[
  (R,e_R,\alpha,\beta):
  \mathcal A\Longrightarrow\mathcal A'
\]
induces initialized Mealy simulations
\[
  \operatorname{out}(\mathcal A)
  \Longrightarrow
  \operatorname{out}(\mathcal A')
\]
and
\[
  \operatorname{cl}_{\mathbf K}(\mathcal A)
  \Longrightarrow
  \operatorname{cl}_{\mathbf K}(\mathcal A').
\]
They depend only on the connected component of the raw optic simulation and preserve identities and
composition up to canonical coherence.
\end{proposition}

\begin{proof}
The output comparison is the pasting
\[
\begin{aligned}
  ((!_M)_\#\times1)\mathbf F
  &\Longrightarrow
  ((!_M)_\#\times1)(R\times1)\mathbf F'
  \\
  &\cong
  \bigl(
    ((!_M)_\#R)\times1
  \bigr)\mathbf F'
  \\
  &\overset{\epsilon_R\times1}{\Longrightarrow}
  ((!_N)_\#\times1)\mathbf F'.
\end{aligned}
\]
The closure comparison is
\[
\begin{aligned}
  \mathbf G(1_M\times\mathbf K)\mathbf F
  &\Longrightarrow
  \mathbf G(1_M\times\mathbf K)(R\times1)\mathbf F'
  \\
  &\cong
  \mathbf G(R\times1)(1_N\times\mathbf K)\mathbf F'
  \\
  &\Longrightarrow
  \mathbf G'(1_N\times\mathbf K)\mathbf F'.
\end{aligned}
\]

If \(\rho:R\Longrightarrow R'\) is a modification, discard preservation gives
\[
  \epsilon_R
  =
  \epsilon_{R'}
  \circ
  ((!_M)_\#\ast\rho),
\]
so the output pastings agree. The closure pastings agree by the defining modification equations.
\end{proof}

\section{Residual coidentification}
\label{sec:residual-coidentification}

Ordinary optic normal forms identify presentations along pure residual slides. A noninvertible
residual slide cannot be used as an invertible bicategorical compositor. It must instead be made an
identity by a homwise categorical coidentifier.

\begin{definition}[Pure residual slide]
\label{def:pure-residual-slide}
Let
\[
  \mathcal A=(M,\mathbf F,\mathbf G),
  \qquad
  \mathcal A'=(N,\mathbf F',\mathbf G')
\]
be initialized optic presentations.

A pure residual slide from \(\mathcal A\) to \(\mathcal A'\) consists of a map
\[
  h:N\to M
\]
and strict invertible initialized simulations
\[
  \phi:
  \mathbf F
  \overset{\cong}{\Longrightarrow}
  \bigl(
    \operatorname{Tot}(h_\#)\times1_{B^+}
  \bigr)
  \circ\mathbf F',
\]
\[
  \psi:
  \mathbf G\circ
  \bigl(
    \operatorname{Tot}(h_\#)\times1_{B^-}
  \bigr)
  \overset{\cong}{\Longrightarrow}
  \mathbf G'.
\]
The residual reparametrizer is \(h_\#\), equipped with its canonical discard section.

When the initialized leg data have been chosen so that the two displayed comparisons are literal
equalities, \(\phi\) and \(\psi\) are identity simulations.
\end{definition}

Composite pure slides are defined using the canonical map-representable comparison
\[
  h_\#\circ k_\#
  \cong
  (h\circ k)_\#.
\]
Tensor products use
\[
  h_\#\times k_\#
  \cong
  (h\times k)_\#.
\]
Since the leg cells are strict and invertible, pure residual slides are closed under composition,
whiskering, and tensor.

For signed interfaces \(A,B\), put
\[
  \mathcal H_{A,B}
  :=
  \operatorname{MOpt}_{\mathrm{op}}(\mathcal E)(A,B).
\]
Let \(\mathcal P_{A,B}\) be the discrete category whose objects are pure residual slides in
\(\mathcal H_{A,B}\). Their sources and targets define functors
\[
  s,t:
  \mathcal P_{A,B}
  \rightrightarrows
  \mathcal H_{A,B},
\]
and the slides themselves define a natural transformation
\[
  \eta:s\Longrightarrow t.
\]

\begin{definition}[Residual coidentifier]
\label{def:residual-coidentifier}
The residual-codescent hom-category is the coidentifier
\[
  q_{A,B}:
  \mathcal H_{A,B}
  \longrightarrow
  \overline{\mathcal H}_{A,B}
\]
of \(\eta\) in \(\mathbf{Cat}\). Thus
\[
  q_{A,B}s=q_{A,B}t
\]
and
\[
  q_{A,B}\eta=1.
\]
Equivalently, \(\overline{\mathcal H}_{A,B}\) is the quotient by the least category congruence that
identifies the source and target of every pure residual slide and sends the slide itself to the
identity of the resulting object class.
\end{definition}

\begin{theorem}[Residual-codescent optic bicategory]
\label{thm:residual-codescent-optics}
The hom-categories
\[
  \overline{\mathcal H}_{A,B}
\]
form a symmetric monoidal bicategory
\[
  \overline{\operatorname{MOpt}}_{\mathrm{op}}(\mathcal E).
\]
The quotient pseudofunctor
\[
  Q_{\mathrm{res}}:
  \operatorname{MOpt}_{\mathrm{op}}(\mathcal E)
  \longrightarrow
  \overline{\operatorname{MOpt}}_{\mathrm{op}}(\mathcal E)
\]
is symmetric monoidal.
\end{theorem}

\begin{proof}
Whiskering a pure residual slide by an optic replaces its residual map by a product with an identity
and pastes the canonical strict leg comparisons. It is therefore again a pure residual slide.
Horizontal composition and tensor of pure residual slides are represented by products of their
residual maps and are again pure.

Consequently horizontal composition and tensor coidentify the generating transformations. By the
universal property of the homwise coidentifiers, they descend uniquely to the quotient
hom-categories. The associators, unitors, and symmetry descend from
\(\operatorname{MOpt}_{\mathrm{op}}(\mathcal E)\), and their coherence equations remain valid.
\end{proof}

\begin{definition}[Raw tight optic simulation]
\label{def:raw-tight-optic-simulation}
A raw operational optic simulation is tight when its residual reparametrizer is map-representable
and its two leg simulations are strict initialized simulations.
\end{definition}

For each \(A,B\), define
\[
  \overline{\mathcal H}^{\mathrm{tight}}_{A,B}
  \subseteq
  \overline{\mathcal H}_{A,B}
\]
to be the wide subcategory generated by the images under \(q_{A,B}\) of raw tight optic
simulations.

\begin{proposition}[Tight locally wide fragment]
\label{prop:tight-fragment}
The hom-categories
\[
  \overline{\mathcal H}^{\mathrm{tight}}_{A,B}
\]
form a symmetric monoidal locally wide sub-bicategory
\[
  \overline{\operatorname{MOpt}}_{\mathrm{tight}}(\mathcal E)
  \subseteq
  \overline{\operatorname{MOpt}}_{\mathrm{op}}(\mathcal E)
\]
with the same objects and \(1\)-cells.
\end{proposition}

\begin{proof}
Composite and tensor products of map-representable residual changes are map-representable.
Horizontal and vertical composites of strict initialized simulations are strict. Pure slides are
themselves generated by map-representable strict cells and become identities after codescent.

Horizontal composition and tensor therefore preserve the generating tight arrows and hence the wide
subcategories they generate. The coherence optics use finite-product coherence functors and thus
belong to the same generated fragment.
\end{proof}

The full operational optic bicategory retains stateful residual reparametrizers. Residual
coidentification kills only the pure map-representable slides used to change residual presentation.

\section{Frozen parameters and parametrized arenas}
\label{sec:parametrized-arenas}

Let
\[
  \operatorname{Disc}:\mathcal E\longrightarrow\operatorname{Cat}(\mathcal E)
\]
be the discrete internal-category functor. For \(\Sigma\in\mathcal E\), define the signed parameter
interface
\[
  \widehat\Sigma
  :=
  (\operatorname{Disc}(\Sigma),\mathbbm1).
\]
Then
\[
  \widehat\Sigma\otimes(A^+,A^-)
  \cong
  (\operatorname{Disc}(\Sigma)\times A^+,A^-).
\]

For a map
\[
  u:T\longrightarrow\Sigma,
\]
write
\[
  \widehat u:
  \widehat T\longrightarrow\widehat\Sigma
\]
for the signed map-representable optic induced by
\(\operatorname{Disc}(u)\) on the positive component and the identity of \(\mathbbm1\) on the
negative component.

\begin{definition}[Parametrized Mealy arena]
\label{def:parametrized-mealy-arena}
A parametrized Mealy arena
\[
  (\Sigma,\mathcal A):A\rightsquigarrow B
\]
consists of an object \(\Sigma\in\mathcal E\) and an operational Mealy optic
\[
  \mathcal A:
  \widehat\Sigma\otimes A
  \rightsquigarrow_{\mathrm{opt}}
  B.
\]
A representative consists of a residual category \(M\) and initialized Mealy morphisms
\[
  \mathbf F:
  \operatorname{Disc}(\Sigma)\times A^+
  \rightsquigarrow
  M\times B^+,
\]
\[
  \mathbf G:
  M\times B^-
  \rightsquigarrow
  A^-.
\]
\end{definition}

The parameter is frozen because its interface is discrete. Persistent adaptation is carried by the
Mealy state of the legs, not by arrows of \(\operatorname{Disc}(\Sigma)\).

\begin{definition}[Parametrized arena simulation]
\label{def:parametrized-arena-simulation}
Let
\[
  (\Sigma,\mathcal A),
  \qquad
  (T,\mathcal B)
\]
have the same signed source and target. A parametrized arena simulation
\[
  (u,\Theta):
  (\Sigma,\mathcal A)
  \Longrightarrow
  (T,\mathcal B)
\]
consists of a map
\[
  u:T\longrightarrow\Sigma
\]
and an operational optic simulation
\[
  \Theta:
  \mathcal A\circ_{\mathrm{opt}}(\widehat u\otimes1_A)
  \Longrightarrow
  \mathcal B.
\]
\end{definition}

The parameter direction is contravariant. An implementation indexed by \(T\) is represented in the
parameterization \(\Sigma\) by \(u:T\to\Sigma\).

For
\[
  (\Sigma,\mathcal A):A\rightsquigarrow B,
  \qquad
  (T,\mathcal B):B\rightsquigarrow C,
\]
the composite has parameter object \(\Sigma\times T\) and arena
\[
\begin{aligned}
  \widehat{\Sigma\times T}\otimes A
  &\cong
  \widehat T\otimes(\widehat\Sigma\otimes A)
  \\
  &\xrightsquigarrow{1_{\widehat T}\otimes\mathcal A}
  \widehat T\otimes B
  \xrightsquigarrow{\mathcal B}
  C.
\end{aligned}
\]

\begin{theorem}[Parametrized operational arenas]
\label{thm:parametrized-operational-arenas}
Parametrized operational Mealy arenas and parametrized arena simulations form a symmetric monoidal
bicategory
\[
  \operatorname{ParaMOpt}_{\mathrm{op}}(\mathcal E).
\]
Residual coidentification gives a symmetric monoidal bicategory
\[
  \overline{\operatorname{ParaMOpt}}_{\mathrm{op}}(\mathcal E),
\]
and the tight locally wide fragment is closed under parameterization, composition, and tensor.
\end{theorem}

\begin{proof}
This is the parameter construction for the cartesian action of \(\mathcal E\) on
\(\operatorname{MOpt}_{\mathrm{op}}(\mathcal E)\) through
\(\Sigma\mapsto\widehat\Sigma\).

Vertical composition is
\[
  (v,\Xi)\circ(u,\Theta)
  =
  \bigl(
    u\circ v,
    \Xi\circ(\Theta\ast(\widehat v\otimes1))
  \bigr).
\]
Horizontal composition multiplies parameter maps and pastes optic simulations. Tensor uses
cartesian product of parameter objects.

Pure residual slides are stable under the parameter action, so the construction descends through
the homwise coidentifiers. Tightness is preserved by the same operations.
\end{proof}

\subsection{Generalized incumbents}

A generalized incumbent is a map
\[
  \sigma:U\longrightarrow\Sigma.
\]
All internally defined strategic constructions may be pulled back along \(\sigma\) to the slice
topos \(\mathcal E/U\). Equivalently, one may precompose the arena by
\[
  \widehat\sigma\otimes1_A.
\]
Thus no global-point condition is imposed. Statements made using an internal variable
\(\sigma:\Sigma\) are stable under arbitrary generalized incumbents.

\section{Internal strategic context doctrines}
\label{sec:strategic-context-doctrines}

A strategic context doctrine specifies which environmental data are presented to the response
machine and how an outer context decomposes under sequential and parallel composition. The doctrine
is internal to \(\mathcal E\).

Let \(\mathcal K\) be a symmetric monoidal bicategory of parametrized arenas.

\begin{definition}[Internal strategic context doctrine]
\label{def:strategic-context-doctrine}
An internal strategic context doctrine
\[
  \mathfrak C
\]
on \(\mathcal K\) consists of the following data.

For every arena
\[
  \mathcal A:A\rightsquigarrow_\Sigma B,
\]
there is an internal category
\[
  \mathfrak C(\mathcal A).
\]

For every parametrized arena simulation
\[
  \gamma=(u,\Theta):
  \mathcal A\Longrightarrow\mathcal A',
\]
there is an internal functor
\[
  \gamma^*:
  \mathfrak C(\mathcal A')
  \longrightarrow
  \mathfrak C(\mathcal A),
\]
and these functors form, homwise, a contravariant pseudofunctor
\[
  \mathfrak C_{A,B}:
  \mathcal K(A,B)^{\mathrm{op}}
  \longrightarrow
  \operatorname{Cat}(\mathcal E).
\]

For composable arenas
\[
  \mathcal A:A\rightsquigarrow_\Sigma B,
  \qquad
  \mathcal B:B\rightsquigarrow_T C,
\]
there is an internal sequential decomposition functor
\[
\begin{aligned}
  \Lambda^{\mathfrak C}_{\mathcal B,\mathcal A}:
  \operatorname{Disc}(\Sigma\times T)
  \times
  \mathfrak C(\mathcal B\circ\mathcal A)
  \longrightarrow
  \mathfrak C(\mathcal A)
  \times
  \mathfrak C(\mathcal B).
\end{aligned}
\]

For arenas
\[
  \mathcal A:A\rightsquigarrow_\Sigma B,
  \qquad
  \mathcal H:C\rightsquigarrow_T D,
\]
there is an internal parallel decomposition functor
\[
\begin{aligned}
  \Pi^{\mathfrak C}_{\mathcal A,\mathcal H}:
  \operatorname{Disc}(\Sigma\times T)
  \times
  \mathfrak C(\mathcal A\otimes\mathcal H)
  \longrightarrow
  \mathfrak C(\mathcal A)
  \times
  \mathfrak C(\mathcal H).
\end{aligned}
\]

For arena simulations
\[
  \gamma=(u,\Theta):
  \mathcal A\Longrightarrow\mathcal A',
  \qquad
  \delta=(v,\Xi):
  \mathcal B\Longrightarrow\mathcal B',
\]
there is an invertible internal natural transformation
\[
\begin{aligned}
  n^{\Lambda}_{\delta,\gamma}:&
  (\gamma^*\times\delta^*)
  \circ
  \Lambda^{\mathfrak C}_{\mathcal B',\mathcal A'}
  \\
  &\Longrightarrow
  \Lambda^{\mathfrak C}_{\mathcal B,\mathcal A}
  \circ
  \bigl(
    \operatorname{Disc}(u\times v)
    \times
    (\delta\circ\gamma)^*
  \bigr).
\end{aligned}
\]
There is an analogous invertible natural transformation
\[
  n^{\Pi}_{\gamma,\eta}
\]
for parallel decomposition.

For three composable arenas
\[
  \mathcal A:A\rightsquigarrow_\Sigma B,
  \qquad
  \mathcal B:B\rightsquigarrow_T C,
  \qquad
  \mathcal L:C\rightsquigarrow_U D,
\]
let
\[
  \Lambda^{L}_{\mathcal L,\mathcal B,\mathcal A},
  \qquad
  \Lambda^{R}_{\mathcal L,\mathcal B,\mathcal A}
\]
be the two functors
\[
\begin{aligned}
  \operatorname{Disc}(\Sigma\times T\times U)
  \times
  \mathfrak C(\mathcal L\circ\mathcal B\circ\mathcal A)
  \longrightarrow
  \mathfrak C(\mathcal A)
  \times
  \mathfrak C(\mathcal B)
  \times
  \mathfrak C(\mathcal L)
\end{aligned}
\]
obtained respectively by first decomposing
\[
  \mathcal L\circ(\mathcal B\circ\mathcal A)
\]
and by first decomposing
\[
  (\mathcal L\circ\mathcal B)\circ\mathcal A,
\]
using the projections from \(\Sigma\times T\times U\) to supply the local strategy coordinates.

There is an invertible internal natural transformation
\[
  a^{\Lambda}_{\mathcal L,\mathcal B,\mathcal A}:
  \Lambda^{L}_{\mathcal L,\mathcal B,\mathcal A}
  \Longrightarrow
  \Lambda^{R}_{\mathcal L,\mathcal B,\mathcal A}.
\]

There are similarly typed invertible internal natural transformations for:

\begin{enumerate}
  \item left and right sequential units;
  \item parallel associativity;
  \item parallel left and right units;
  \item parallel symmetry;
  \item sequential--parallel middle-four interchange.
\end{enumerate}

The middle-four comparison has source and target the two functors
\[
\begin{aligned}
&
\operatorname{Disc}
\bigl(
  \Sigma_1\times\Sigma_2\times T_1\times T_2
\bigr)
\times
\mathfrak C
\bigl(
  (\mathcal B_1\otimes\mathcal B_2)
  \circ
  (\mathcal A_1\otimes\mathcal A_2)
\bigr)
\\
&\hspace{20mm}\longrightarrow
\mathfrak C(\mathcal A_1)
\times
\mathfrak C(\mathcal A_2)
\times
\mathfrak C(\mathcal B_1)
\times
\mathfrak C(\mathcal B_2)
\end{aligned}
\]
obtained in the two orders:

\begin{enumerate}
  \item sequential decomposition followed by parallel decomposition of each local context;
  \item parallel decomposition followed by sequential decomposition of each tensor factor.
\end{enumerate}

These transformations satisfy the triangle, pentagon, hexagon, middle-four interchange, and
pseudonaturality equations.
\end{definition}

The explicit pseudonaturality and coherence transformations are part of the doctrine. They are the
comparison cells needed when evaluator machines are reindexed during composition of strategic
simulations.

\section{Internal counterfactual query machines}
\label{sec:counterfactual-query-machines}

\subsection{Internal cores}

Let \(\mathcal C\) be an internal category. Define
\[
  \operatorname{Inv}(\mathcal C)
\]
to be the object of pairs \((f,g)\) of arrows satisfying
\[
  g\circ f=1,
  \qquad
  f\circ g=1.
\]
Projection to \(f\) is monic because inverses are unique.

\begin{definition}[Internal core]
\label{def:internal-core}
The internal core
\[
  \mathcal C^\simeq
\]
has object object \(\mathcal C_0\) and arrow object the image of
\[
  \operatorname{Inv}(\mathcal C)\longrightarrow\mathcal C_1.
\]
It is an internal groupoid.
\end{definition}

\begin{lemma}[Functoriality of the core]
\label{lem:core-functoriality}
Every internal functor
\[
  F:\mathcal C\to\mathcal D
\]
restricts to an internal functor
\[
  F^\simeq:\mathcal C^\simeq\to\mathcal D^\simeq.
\]
Every invertible internal natural transformation restricts to an invertible natural transformation
between the induced core functors.
\end{lemma}

\begin{proof}
A functor sends inverse pairs to inverse pairs. Naturality and invertibility are inherited from the
ambient internal categories.
\end{proof}

\subsection{Internal deviation categories}

Let
\[
  \Sigma^\ast
  :=
  \coprod_{n\in\mathbb N}\Sigma^n
\]
be the internal free monoid on \(\Sigma\), with unit the empty word and multiplication given by list
concatenation.

\begin{definition}[Internal deviation category]
\label{def:internal-deviation-category}
The internal deviation category
\[
  \mathsf{Dev}(\Sigma)
\]
has
\[
  \mathsf{Dev}(\Sigma)_0=\Sigma,
\]
\[
  \mathsf{Dev}(\Sigma)_1=\Sigma\times\Sigma^\ast,
\]
with both source and target equal to the first projection. Identity is the empty word and
composition is concatenation.

The one-letter deviation associated to an alternative \(\sigma'\) is written
\[
  [\sigma']:\sigma\longrightarrow\sigma.
\]
\end{definition}

A map
\[
  u:T\to\Sigma
\]
induces an internal functor
\[
  \mathsf{Dev}(u):
  \mathsf{Dev}(T)
  \longrightarrow
  \mathsf{Dev}(\Sigma)
\]
by applying \(u\) to the incumbent and to every letter of the deviation word.

\begin{definition}[Internal query category]
\label{def:internal-query-category}
For an internal category \(\mathcal C\), define
\[
  \mathsf{Qry}(\Sigma,\mathcal C)
  :=
  \mathsf{Dev}(\Sigma)\times\mathcal C^\simeq.
\]
Its objects are pairs
\[
  (\sigma,c).
\]
A one-query deviation is
\[
  ([\sigma'],1_c):
  (\sigma,c)\longrightarrow(\sigma,c),
\]
and a context isomorphism
\[
  f:c\overset{\cong}{\longrightarrow}d
\]
gives
\[
  (1_\sigma,f):
  (\sigma,c)\longrightarrow(\sigma,d).
\]
\end{definition}

A map \(u:T\to\Sigma\) and an internal functor \(F:\mathcal D\to\mathcal C\) induce
\[
  \mathsf{Qry}(u,F):
  \mathsf{Qry}(T,\mathcal D)
  \longrightarrow
  \mathsf{Qry}(\Sigma,\mathcal C).
\]
An invertible natural transformation
\[
  \eta:F\Longrightarrow G
\]
induces an invertible natural transformation
\[
  \mathsf{Qry}(u,\eta):
  \mathsf{Qry}(u,F)
  \Longrightarrow
  \mathsf{Qry}(u,G).
\]

\subsection{Truth output}

Let \(\Omega\) be the subobject classifier of \(\mathcal E\).

\begin{definition}[Internal truth category]
\label{def:internal-truth-category}
The internal truth category
\[
  \mathsf{Truth}_\wedge
\]
has one object,
\[
  (\mathsf{Truth}_\wedge)_0=\mathbbm1,
\]
arrow object
\[
  (\mathsf{Truth}_\wedge)_1=\Omega,
\]
identity
\[
  \top:\mathbbm1\to\Omega,
\]
and composition
\[
  \wedge:\Omega\times\Omega\to\Omega.
\]
\end{definition}

Meet defines a strict symmetric monoidal internal functor
\[
  \wedge:
  \mathsf{Truth}_\wedge\times\mathsf{Truth}_\wedge
  \longrightarrow
  \mathsf{Truth}_\wedge.
\]

\subsection{Initialized response evaluators}

\begin{definition}[Counterfactual response evaluator]
\label{def:counterfactual-response-evaluator}
A counterfactual response evaluator for \((\Sigma,\mathcal C)\) consists of:
\begin{enumerate}
  \item a transition-closed initialized Mealy morphism
  \[
    \mathbf E:
    \mathsf{Qry}(\Sigma,\mathcal C)
    \rightsquigarrow
    \mathsf{Truth}_\wedge;
  \]
  \item a ready-state map
  \[
    \rho_E:
    \Sigma\times\mathcal C_0
    \longrightarrow
    J_E
  \]
  over the query-object map;
  \item context transparency: for every context isomorphism
  \[
    f:c\overset{\cong}{\longrightarrow}d
  \]
  and every admissible state \(i\) over \((\sigma,d)\),
  \[
    \omega_E
    \bigl(
      (1_\sigma,f),
      j_Ei
    \bigr)
    =
    \top,
  \]
  and
  \[
    \lambda_E
    \bigl(
      (1_\sigma,f),
      \rho_E(\sigma,d)
    \bigr)
    =
    \rho_E(\sigma,c).
  \]
\end{enumerate}
\end{definition}

Context transparency is imposed on every admissible initialized state, not only on ready states.
Since \(f^{-1}\) is also a context isomorphism, the action laws make context transport invertible.

\begin{definition}[One-query classifier]
\label{def:one-query-classifier}
The one-query classifier of \(E\) is the morphism
\[
  \chi_E:
  \Sigma\times\mathcal C_0\times\Sigma
  \longrightarrow
  \Omega
\]
obtained by reading the output produced when the ready state
\(\rho_E(\sigma,c)\) processes the one-letter deviation
\[
  ([\sigma'],1_c).
\]

It classifies the best-response subobject
\[
  \mathsf{BR}_E
  \hookrightarrow
  \Sigma\times\mathcal C_0\times\Sigma.
\]
The equilibrium subobject is the pullback of \(\mathsf{BR}_E\) along
\[
  \Sigma\times\mathcal C_0
  \longrightarrow
  \Sigma\times\mathcal C_0\times\Sigma,
  \qquad
  (\sigma,c)\longmapsto(\sigma,c,\sigma).
\]
\end{definition}

\begin{proposition}[Context invariance]
\label{prop:context-invariance}
For every context isomorphism
\[
  f:c\overset{\cong}{\longrightarrow}d,
\]
one has internally
\[
  \chi_E(\sigma,c,\sigma')
  =
  \chi_E(\sigma,d,\sigma').
\]
\end{proposition}

\begin{proof}
In the product category
\(\mathsf{Dev}(\Sigma)\times\mathcal C^\simeq\), deviation and context transport commute:
\[
  (1_\sigma,f)\circ([\sigma'],1_c)
  =
  ([\sigma'],1_d)\circ(1_\sigma,f).
\]
Evaluate the two factorizations from the ready state over \(d\). The context arrow emits
\(\top\) before or after the deviation. The Mealy multiplication law therefore identifies the two
deviation outputs. Ready-state transport identifies the state over \(c\) with
\(\rho_E(\sigma,c)\).
\end{proof}

\subsection{Extensional evaluators}

Let
\[
  \operatorname{fold}_\wedge:
  \Omega^\ast\longrightarrow\Omega
\]
be finite meet, with empty meet \(\top\).

\begin{definition}[Extensional response evaluator]
\label{def:extensional-response-evaluator}
Let
\[
  \chi:
  \Sigma\times\mathcal C_0\times\Sigma
  \longrightarrow
  \Omega
\]
be invariant under arrows of \(\mathcal C^\simeq\).

Define an internal functor
\[
  \mathsf{Qry}(\chi):
  \mathsf{Qry}(\Sigma,\mathcal C)
  \longrightarrow
  \mathsf{Truth}_\wedge
\]
by sending a query arrow consisting of a context isomorphism and a deviation word
\[
  \sigma'_1\cdots\sigma'_n
\]
to
\[
  \bigwedge_{i=1}^{n}
  \chi(\sigma,c,\sigma'_i).
\]
Context invariance makes this independent of whether \(c\) is read at the source or target of the
context isomorphism.

The associated totally initialized map-representable Mealy morphism is the extensional evaluator
classified by \(\chi\).
\end{definition}

\begin{proposition}[Canonical extensional realization]
\label{prop:canonical-extensional-realization}
Every context-isomorphism-invariant subobject
\[
  \mathsf{BR}
  \hookrightarrow
  \Sigma\times\mathcal C_0\times\Sigma
\]
has a canonical extensional response evaluator.
\end{proposition}

\begin{proof}
Take the classifying morphism
\[
  \chi_{\mathsf{BR}}:
  \Sigma\times\mathcal C_0\times\Sigma
  \to\Omega
\]
and apply Definition~\ref{def:extensional-response-evaluator}.
\end{proof}

\subsection{Transport along context isomorphisms}

\begin{lemma}[Evaluator transport]
\label{lem:evaluator-transport}
Let \(E\) be a context-transparent evaluator for \((\Sigma,\mathcal C)\). Let
\[
  u:T\to\Sigma,
\]
let
\[
  F,G:\mathcal D\to\mathcal C
\]
be internal functors, and let
\[
  \eta:F\Longrightarrow G
\]
be an invertible internal natural transformation. Then \(\eta\) induces an invertible
response-strict initialized Mealy simulation
\[
  \operatorname{tr}_E(u,\eta):
  E\circ\mathsf{Qry}(u,F)_\#
  \Longrightarrow
  E\circ\mathsf{Qry}(u,G)_\#.
\]
\end{lemma}

\begin{proof}
A state of the target composite over \((\tau,d)\) is an admissible \(E\)-state over
\[
  (u\tau,Gd).
\]
Process the context-isomorphism arrow
\[
  (1_{u\tau},\eta_d):
  (u\tau,Fd)\longrightarrow(u\tau,Gd)
\]
at that state. Target-to-source processing produces a state over \((u\tau,Fd)\), defining the
carrier comparison.

Naturality of \(\eta\) gives compatibility with every query arrow. Context transparency says that
the processed comparison arrow emits \(\top\), so the output comparison is the identity truth
arrow. The ready-state equation follows from ready-state transport. Processing
\(\eta^{-1}\) supplies the inverse.
\end{proof}

\subsection{Response-strict evaluator simulations}

\begin{definition}[Evaluator simulation over reindexing]
\label{def:evaluator-simulation-over-reindexing}
Let \(E\) be an evaluator for \((\Sigma,\mathcal C)\), and let \(E'\) be an evaluator for
\((T,\mathcal D)\).

An evaluator simulation over
\[
  (u,F),
  \qquad
  u:T\to\Sigma,
  \quad
  F:\mathcal D\to\mathcal C,
\]
is an initialized Mealy simulation
\[
  E\circ\mathsf{Qry}(u,F)_\#
  \Longrightarrow
  E'
\]
whose carrier comparison preserves ready states.

It is response-strict when its output-comparison component is everywhere \(\top\).
\end{definition}

\begin{proposition}[Classifier preservation and response-strict closure]
\label{prop:classifier-preservation}
A response-strict evaluator simulation over \((u,F)\) satisfies
\[
  \chi_{E'}(\tau,d,\tau')
  =
  \chi_E(u\tau,Fd,u\tau').
\]

Response-strict evaluator simulations are closed under vertical composition, left and right
whiskering by fixed initialized Mealy \(1\)-cells, product tensor, evaluator transport, and
postcomposition by the meet functor.
\end{proposition}

\begin{proof}
Apply the Mealy output equation to a one-query generator at a preserved ready state. Since the output
comparison is \(\top\), the two emitted truth arrows are equal.

Vertical composition and tensor use
\[
  \top\wedge\top=\top.
\]
For horizontal composition, the downstream state is transported along the upstream
output-comparison arrow. When the upstream comparison is strict, this arrow is an identity. The
Mealy unit law says that processing it emits an identity output. Composing with the strict
downstream output comparison again gives an identity. Hence strict simulations are closed under
horizontal composition and therefore under ordinary bicategorical whiskering.

Evaluator transport is strict by context transparency, and the meet functor sends the pair of
identity truth arrows to \(\top\).
\end{proof}

\section{The bicategory of strategic Mealy games}
\label{sec:strategic-mealy-games}

Let \(\mathcal K\) be a symmetric monoidal bicategory of parametrized arenas and let
\(\mathfrak C\) be an internal strategic context doctrine on \(\mathcal K\).

\begin{definition}[Strategic Mealy game]
\label{def:strategic-mealy-game}
A strategic Mealy game
\[
  \mathcal G:A\longrightarrow B
\]
relative to \(\mathfrak C\) consists of:
\begin{enumerate}
  \item a strategy object
  \[
    \Sigma_{\mathcal G}\in\mathcal E;
  \]
  \item a parametrized arena
  \[
    \mathcal A_{\mathcal G}:
    A\rightsquigarrow_{\Sigma_{\mathcal G}}B
  \]
  in \(\mathcal K\);
  \item a context-transparent counterfactual evaluator
  \[
    E_{\mathcal G}:
    \mathsf{Qry}
    \bigl(
      \Sigma_{\mathcal G},
      \mathfrak C(\mathcal A_{\mathcal G})
    \bigr)
    \rightsquigarrow
    \mathsf{Truth}_\wedge.
  \]
\end{enumerate}
\end{definition}

The extensional strategic shadow is the arena together with the classifier
\[
  \chi_{\mathcal G}:
  \Sigma_{\mathcal G}
  \times
  \mathfrak C(\mathcal A_{\mathcal G})_0
  \times
  \Sigma_{\mathcal G}
  \longrightarrow
  \Omega.
\]
A stateful evaluator may contain histories, cached simulations, beliefs, or adaptive query state
that are invisible in this one-query shadow.

\subsection{Sequential composition}

Let
\[
  \mathcal G:A\to B,
  \qquad
  \mathcal H:B\to C
\]
have strategy objects \(\Sigma,T\). The sequential context decomposition functor induces an internal
query-decomposition functor
\[
\begin{aligned}
  D_{\mathcal H,\mathcal G}:&
  \mathsf{Qry}
  \bigl(
    \Sigma\times T,
    \mathfrak C(
      \mathcal A_{\mathcal H}\circ\mathcal A_{\mathcal G}
    )
  \bigr)
  \\
  &\longrightarrow
  \mathsf{Qry}
  \bigl(
    \Sigma,
    \mathfrak C(\mathcal A_{\mathcal G})
  \bigr)
  \times
  \mathsf{Qry}
  \bigl(
    T,
    \mathfrak C(\mathcal A_{\mathcal H})
  \bigr).
\end{aligned}
\]
On objects it is
\[
  ((\sigma,\tau),c)
  \longmapsto
  \bigl(
    (\sigma,c_{\mathcal G}),
    (\tau,c_{\mathcal H})
  \bigr),
\]
where
\[
  (c_{\mathcal G},c_{\mathcal H})
  =
  \Lambda^{\mathfrak C}_{\mathcal H,\mathcal G}
  ((\sigma,\tau),c).
\]
On a deviation word in \(\Sigma\times T\), it applies the two coordinate projections to every
letter. On context isomorphisms, it uses the functorial action of
\(\Lambda^{\mathfrak C}_{\mathcal H,\mathcal G}\).

\begin{definition}[Sequential composite]
\label{def:strategic-sequential-composite}
The composite
\[
  \mathcal H\circ\mathcal G
\]
has strategy object
\[
  \Sigma\times T,
\]
arena
\[
  \mathcal A_{\mathcal H}\circ\mathcal A_{\mathcal G},
\]
and evaluator
\[
  E_{\mathcal H\circ\mathcal G}
  :=
  \wedge_\#
  \circ
  (E_{\mathcal G}\times E_{\mathcal H})
  \circ
  (D_{\mathcal H,\mathcal G})_\#.
\]
Its ready-state map is
\[
\begin{aligned}
  \rho_{\mathcal H\circ\mathcal G}
  ((\sigma,\tau),c)
  :=
  \bigl(
    \rho_{\mathcal G}(\sigma,c_{\mathcal G}),
    \rho_{\mathcal H}(\tau,c_{\mathcal H})
  \bigr),
\end{aligned}
\]
under the canonical carrier and initialization identifications for product and pipeline
composition.
\end{definition}

Its one-query classifier is
\[
\begin{aligned}
&\chi_{\mathcal H\circ\mathcal G}
\bigl(
  (\sigma,\tau),
  c,
  (\sigma',\tau')
\bigr)
\\
&\qquad=
\chi_{\mathcal G}
(\sigma,c_{\mathcal G},\sigma')
\wedge
\chi_{\mathcal H}
(\tau,c_{\mathcal H},\tau').
\end{aligned}
\]

The identity game has strategy object \(\mathbbm1\), identity arena, and constantly true evaluator.
Tensor uses parallel context decomposition and the same meet functor.

\subsection{Strategic simulations}

\begin{definition}[Response-strict strategic simulation]
\label{def:strategic-game-simulation}
Let
\[
  \mathcal G,\mathcal H:A\to B
\]
have strategy objects \(\Sigma,T\).

A response-strict strategic simulation
\[
  (u,\Theta,\Xi):
  \mathcal G\Longrightarrow\mathcal H
\]
consists of:
\begin{enumerate}
  \item a map
  \[
    u:T\to\Sigma;
  \]
  \item a parametrized arena simulation
  \[
    (u,\Theta):
    (\Sigma,\mathcal A_{\mathcal G})
    \Longrightarrow
    (T,\mathcal A_{\mathcal H});
  \]
  \item a response-strict evaluator simulation
  \[
    \Xi:
    E_{\mathcal G}
    \circ
    \mathsf{Qry}
    \bigl(
      u,
      (u,\Theta)^*
    \bigr)_\#
    \Longrightarrow
    E_{\mathcal H}.
  \]
\end{enumerate}
\end{definition}

Its classifier equation is
\[
  \chi_{\mathcal H}(\tau,c,\tau')
  =
  \chi_{\mathcal G}
  \bigl(
    u\tau,
    (u,\Theta)^*c,
    u\tau'
  \bigr).
\]
Diagonal response and equilibrium are therefore preserved exactly.

\begin{theorem}[Strategic Mealy bicategory]
\label{thm:strategic-mealy-bicategory}
Strategic Mealy games and response-strict strategic simulations form a bicategory
\[
  \operatorname{StrMlyGame}^{=}_{\mathfrak C}(\mathcal K).
\]
If \(\mathfrak C\) carries coherent parallel decomposition, this bicategory is symmetric monoidal.
\end{theorem}

\begin{proof}
The composite evaluator is initialized because \(D_{\mathcal H,\mathcal G}\) is an internal functor,
product tensors initialized evaluators, and meet is an internal functor. Its ready-state map is the
map displayed in Definition~\ref{def:strategic-sequential-composite}. Context transparency is
preserved because the local evaluators emit \(\top\) on context transports and
\[
  \top\wedge\top=\top.
\]

For three games, the doctrine associator
\[
  a^\Lambda
\]
is an invertible internal natural transformation between the two functors assigning the triple of
local contexts. Applying the internal core and \(\mathsf{Qry}\) gives an invertible natural
transformation between the corresponding query decomposition functors.
Lemma~\ref{lem:evaluator-transport} converts it into an invertible response-strict evaluator
simulation. Product coherence and associativity of meet then construct the strategic associator.

The strategic pentagon is the image of the doctrine pentagon under the core construction, the query
construction, evaluator transport, product of evaluators, and postcomposition by meet.

The unitors are obtained from doctrine unit coherence and the unit law for meet.

Vertical composition composes parameter maps, arena simulations, context reindexing functors, and
response-strict evaluator simulations. The pseudofunctorial comparison for context reindexing is
converted by Lemma~\ref{lem:evaluator-transport} into the required evaluator transport.

Horizontal composition uses the explicit doctrine pseudonaturality transformation
\[
  n^\Lambda_{\delta,\gamma},
\]
followed by product and pasting of the local evaluator simulations. Response strictness follows from
Proposition~\ref{prop:classifier-preservation}.

Interchange follows from interchange in the arena and initialized-Mealy bicategories. Parallel
decomposition, doctrine symmetry, and commutativity of meet supply the symmetric monoidal structure.
The triangle, pentagon, hexagon, and middle-four equations reduce to the corresponding doctrine
equations, functoriality of evaluator transport, and coherence of product and meet.
\end{proof}

\section{Tight stateless optics and internal deterministic lenses}
\label{sec:tight-stateless-optics}

For objects \(X,S,Y,R\in\mathcal E\), an internal deterministic lens
\[
  (X,S)\longrightarrow(Y,R)
\]
is a pair
\[
  p:X\to Y,
  \qquad
  c:X\times R\to S.
\]

A tight discrete stateless optic presentation consists of an object \(M\) and maps
\[
  f:X\longrightarrow M\times Y,
  \qquad
  b:M\times R\longrightarrow S,
\]
regarded as totally initialized map-representable Mealy morphisms between discrete internal
categories.

A pure residual slide along
\[
  h:N\to M
\]
is represented by the strict equations
\[
  f=(h\times1_Y)f',
  \qquad
  b(h\times1_R)=b'.
\]

\begin{definition}[Canonical lens presentation]
\label{def:canonical-lens-presentation}
The canonical residual presentation of a lens \((p,c)\) has residual \(X\), forward map
\[
  \bar p(x)=(x,p(x)),
\]
and backward map
\[
  c:X\times R\to S.
\]
\end{definition}

Conversely, a presentation \((M,f,b)\) determines
\[
  p_f:=\pi_Yf
\]
and
\[
  c_{f,b}
  :=
  b\circ(\pi_Mf\times1_R).
\]

\begin{theorem}[Internal lens--optic correspondence]
\label{thm:internal-lens-optic-correspondence}
The locally discrete tight discrete stateless fragment of the residual-codescent optic bicategory is
isomorphic to the locally discrete bicategory associated to internal deterministic lenses.

Homwise, in the fixed universe,
\[
  \operatorname{Lens}_{\mathcal E}((X,S),(Y,R))
  \cong
  \int^{M\in\mathcal E}
  \mathcal E(X,M\times Y)
  \times
  \mathcal E(M\times R,S).
\]
The correspondence preserves identities, sequential composition, and product tensor.
\end{theorem}

\begin{proof}
A pure residual slide does not change the extracted lens. If
\[
  f=(h\times1)f',
  \qquad
  b(h\times1)=b',
\]
then
\[
  \pi_Yf=\pi_Yf'
\]
and
\[
\begin{aligned}
  b(\pi_Mf(x),r)
  &=
  b(h(\pi_Nf'(x)),r)
  \\
  &=
  b'(\pi_Nf'(x),r).
\end{aligned}
\]

Conversely, given \((M,f,b)\), put
\[
  h:=\pi_Mf:X\to M.
\]
The canonical presentation of its extracted lens has forward map
\[
  x\longmapsto(x,\pi_Yf(x))
\]
and backward map
\[
  (x,r)\longmapsto b(h(x),r).
\]
The map \(h\) gives a pure residual slide from the original presentation to the canonical
presentation. Hence every residual-codescent class has a unique canonical representative.

Every raw tight \(2\)-cell between discrete stateless presentations is represented by a map
\[
  h:N\to M
\]
satisfying the two strict equations above. It is therefore a pure residual slide and becomes an
identity in residual codescent. Since the tight quotient hom-category is generated by such cells,
the resulting hom-categories are discrete.

For composable lenses \((p,c)\) and \((q,d)\), raw optic composition has residual \(X\times Y\).
The map
\[
  X\longrightarrow X\times Y,
  \qquad
  x\longmapsto(x,p(x))
\]
gives a pure slide to the canonical residual \(X\). The extracted lens is
\[
  q\circ p,
  \qquad
  (x,u)\longmapsto c(x,d(p(x),u)).
\]
Tensor is coordinatewise.
\end{proof}

\subsection{Parametrized deterministic lenses}

Let
\[
  \overline{\mathcal K}_{\mathrm{td}}
\]
be the symmetric monoidal sub-bicategory of
\[
  \overline{\operatorname{ParaMOpt}}_{\mathrm{op}}(\mathcal E)
\]
whose:
\begin{enumerate}
  \item objects are signed discrete interfaces
  \[
    (\operatorname{Disc}(X),\operatorname{Disc}(S));
  \]
  \item \(1\)-cells have discrete residuals and totally initialized map-representable forward and
  backward legs;
  \item \(2\)-cells belong to the generated tight fragment.
\end{enumerate}

\begin{proposition}[Parametrized lens correspondence]
\label{prop:parametrized-lens-correspondence}
The bicategory \(\overline{\mathcal K}_{\mathrm{td}}\) is symmetrically monoidally isomorphic to the
bicategory of parametrized internal deterministic lenses and strategy reparametrizations.
\end{proposition}

\begin{proof}
Apply Theorem~\ref{thm:internal-lens-optic-correspondence} in each parameter fibre. A parametrized
lens is equivalently a pair
\[
  P:\Sigma\times X\to Y,
  \qquad
  C:\Sigma\times X\times R\to S.
\]
A parameter map
\[
  u:T\to\Sigma
\]
acts by precomposition. For canonical residuals, the strict forward equation forces the residual
map to be
\[
  u\times1_X:T\times X\to\Sigma\times X.
\]
The strict backward equation is exactly preservation of the coplay map. Composition and tensor are
preserved by canonical residualization.
\end{proof}

\section{The stationary context doctrine}
\label{sec:stationary-context-doctrine}

For a signed discrete boundary
\[
  (X,S)\longrightarrow(Y,R),
\]
define the internal stationary context category
\[
  \operatorname{Ctx}_{\mathrm{st}}(X;Y,R)
  :=
  \operatorname{Disc}(X\times R^Y).
\]
An internal element is written
\[
  (x,k),
  \qquad
  k:Y\to R.
\]

Let
\[
  \mathcal G:(X,S)\to(Y,R),
  \qquad
  \mathcal H:(Y,R)\to(Z,Q)
\]
be parametrized deterministic lens processes with strategy objects \(\Sigma,T\), play maps
\(P_{\mathcal G},P_{\mathcal H}\), and coplay maps
\(C_{\mathcal G},C_{\mathcal H}\).

Define
\[
  \kappa_{\mathcal H}:
  T\times Q^Z\times Y
  \longrightarrow
  R
\]
by
\[
  \kappa_{\mathcal H}(\tau,k,y)
  :=
  C_{\mathcal H}
  \bigl(
    \tau,
    y,
    \operatorname{ev}
    (k,P_{\mathcal H}(\tau,y))
  \bigr).
\]
Let
\[
  \bar\kappa_{\mathcal H}:
  T\times Q^Z\longrightarrow R^Y
\]
be its exponential transpose.

The sequential context decomposition is the map
\[
\begin{aligned}
  \Lambda_{\mathcal H,\mathcal G}^{\mathrm{st}}:
  \Sigma\times T\times X\times Q^Z
  &\longrightarrow
  (X\times R^Y)\times(Y\times Q^Z),
  \\
  (\sigma,\tau,x,k)
  &\longmapsto
  \Bigl(
    (x,\bar\kappa_{\mathcal H}(\tau,k)),
    (P_{\mathcal G}(\sigma,x),k)
  \Bigr).
\end{aligned}
\]

For parallel processes
\[
  \mathcal G:(X,S)\to(Y,R),
  \qquad
  \mathcal H:(X',S')\to(Y',R'),
\]
let
\[
  k:(Y\times Y')\to(R\times R').
\]
Define the partially evaluated continuations by exponential transpose of
\[
\begin{aligned}
  &\Sigma\times T\times X'\times(R\times R')^{Y\times Y'}\times Y
  \longrightarrow R,
  \\
  &(\sigma,\tau,x',k,y)
  \longmapsto
  \pi_R
  \operatorname{ev}
  \bigl(
    k,
    (y,P_{\mathcal H}(\tau,x'))
  \bigr),
\end{aligned}
\]
and
\[
\begin{aligned}
  &\Sigma\times T\times X\times(R\times R')^{Y\times Y'}\times Y'
  \longrightarrow R',
  \\
  &(\sigma,\tau,x,k,y')
  \longmapsto
  \pi_{R'}
  \operatorname{ev}
  \bigl(
    k,
    (P_{\mathcal G}(\sigma,x),y')
  \bigr).
\end{aligned}
\]

\begin{proposition}[Stationary context doctrine]
\label{prop:stationary-context-doctrine}
The stationary context categories and the preceding sequential and parallel decomposition maps form
an internal strategic context doctrine
\[
  \mathfrak C_{\mathrm{st}}
\]
on \(\overline{\mathcal K}_{\mathrm{td}}\).
\end{proposition}

\begin{proof}
The context category depends only on the signed boundary. Its reindexing functor along a
boundary-preserving tight arena simulation is therefore the identity.

Pure residual slides do not change the extracted play or coplay maps. Hence both decomposition maps
are invariant under residual coidentification.

For three sequential processes
\[
  \mathcal G,\mathcal H,\mathcal L
\]
with strategy variables \(\sigma,\tau,\upsilon\) and outer continuation \(k\), define
\[
  k_{\mathcal L}(z)
  :=
  C_{\mathcal L}
  \bigl(
    \upsilon,
    z,
    k(P_{\mathcal L}(\upsilon,z))
  \bigr)
\]
and
\[
  k_{\mathcal H,\mathcal L}(y)
  :=
  C_{\mathcal H}
  \bigl(
    \tau,
    y,
    k_{\mathcal L}(P_{\mathcal H}(\tau,y))
  \bigr).
\]

Both sequential parenthesizations assign to \(\mathcal L\) the context
\[
  \bigl(
    P_{\mathcal H}
    (\tau,P_{\mathcal G}(\sigma,x)),
    k
  \bigr),
\]
to \(\mathcal H\) the context
\[
  \bigl(
    P_{\mathcal G}(\sigma,x),
    k_{\mathcal L}
  \bigr),
\]
and to \(\mathcal G\) the context
\[
  (x,k_{\mathcal H,\mathcal L}).
\]
Hence the sequential associativity comparison is equality after canonical cartesian
reassociation.

For a finite tensor and joint continuation \(k\), the local continuation of component \(i\) is
\[
  y_i
  \longmapsto
  \pi_i
  k
  \bigl(
    y_1^\ast,\ldots,y_{i-1}^\ast,
    y_i,
    y_{i+1}^\ast,\ldots,y_n^\ast
  \bigr),
\]
where
\[
  y_j^\ast=P_j(\sigma_j,x_j).
\]
This expression is independent of cartesian parenthesization. Parallel associativity, symmetry, and
middle-four interchange therefore hold by equality after canonical reassociation.

Since all context categories are discrete, every coherence transformation is uniquely determined by
these equalities. The required triangle, pentagon, hexagon, pseudonaturality, and middle-four
equations follow.
\end{proof}

\section{Internal deterministic open games}
\label{sec:internal-open-games}

\begin{definition}[Internal deterministic open game]
\label{def:internal-deterministic-open-game}
An internal deterministic open game
\[
  \mathcal G:(X,S)\longrightarrow(Y,R)
\]
consists of:
\begin{enumerate}
  \item a strategy object
  \[
    \Sigma_{\mathcal G}\in\mathcal E;
  \]
  \item a play map
  \[
    P_{\mathcal G}:
    \Sigma_{\mathcal G}\times X\to Y;
  \]
  \item a coplay map
  \[
    C_{\mathcal G}:
    \Sigma_{\mathcal G}\times X\times R\to S;
  \]
  \item a best-response subobject
  \[
    \mathsf{BR}_{\mathcal G}
    \hookrightarrow
    \Sigma_{\mathcal G}
    \times X
    \times R^Y
    \times
    \Sigma_{\mathcal G}.
  \]
\end{enumerate}
\end{definition}

Write
\[
  \chi_{\mathcal G}:
  \Sigma_{\mathcal G}
  \times X
  \times R^Y
  \times
  \Sigma_{\mathcal G}
  \longrightarrow
  \Omega
\]
for the classifier of \(\mathsf{BR}_{\mathcal G}\).

The equilibrium subobject is classified by
\[
  (\sigma,x,k)
  \longmapsto
  \chi_{\mathcal G}(\sigma,x,k,\sigma).
\]

\begin{definition}[Strategy reparametrization]
\label{def:internal-strategy-reparametrization}
For open games \(\mathcal G,\mathcal H\) with the same boundary, a strategy reparametrization
\[
  u:\mathcal G\Longrightarrow\mathcal H
\]
is a map
\[
  u:\Sigma_{\mathcal H}\longrightarrow\Sigma_{\mathcal G}
\]
such that
\[
  P_{\mathcal H}
  =
  P_{\mathcal G}\circ(u\times1_X),
\]
\[
  C_{\mathcal H}
  =
  C_{\mathcal G}
  \circ
  (u\times1_X\times1_R),
\]
and
\[
  \chi_{\mathcal H}(\tau,x,k,\tau')
  =
  \chi_{\mathcal G}(u\tau,x,k,u\tau').
\]
\end{definition}

\subsection{Sequential composition}

Let
\[
  \mathcal G:(X,S)\to(Y,R),
  \qquad
  \mathcal H:(Y,R)\to(Z,Q)
\]
have strategy objects \(\Sigma,T\). Their composite has strategy object
\[
  \Sigma\times T
\]
and play map
\[
  P_{\mathcal H\circ\mathcal G}
  ((\sigma,\tau),x)
  =
  P_{\mathcal H}
  \bigl(
    \tau,
    P_{\mathcal G}(\sigma,x)
  \bigr).
\]
Its coplay map is
\[
\begin{aligned}
  C_{\mathcal H\circ\mathcal G}
  ((\sigma,\tau),x,q)
  =
  C_{\mathcal G}
  \Bigl(
    \sigma,
    x,
    C_{\mathcal H}
    \bigl(
      \tau,
      P_{\mathcal G}(\sigma,x),
      q
    \bigr)
  \Bigr).
\end{aligned}
\]

Let
\[
  \bar\kappa_{\mathcal H}:
  T\times Q^Z\to R^Y
\]
be the continuation transpose defined in the stationary context doctrine. The composite classifier
is
\[
\begin{aligned}
&\chi_{\mathcal H\circ\mathcal G}
\bigl(
  (\sigma,\tau),
  x,
  k,
  (\sigma',\tau')
\bigr)
\\
&\qquad=
\chi_{\mathcal G}
\bigl(
  \sigma,
  x,
  \bar\kappa_{\mathcal H}(\tau,k),
  \sigma'
\bigr)
\\
&\qquad\qquad\wedge
\chi_{\mathcal H}
\bigl(
  \tau,
  P_{\mathcal G}(\sigma,x),
  k,
  \tau'
\bigr).
\end{aligned}
\]

Tensor has product strategy object, coordinatewise play and coplay, and classifier obtained by meet
after the parallel stationary-context decomposition.

\begin{theorem}[Internal open-game bicategory]
\label{thm:internal-open-game-bicategory}
Signed pairs of objects of \(\mathcal E\), internal deterministic open games, and strategy
reparametrizations form a symmetric monoidal bicategory
\[
  \mathbf{OpenGame}_2(\mathcal E).
\]
Its associators, unitors, and symmetry are the canonical invertible strategy reparametrizations
induced by cartesian product.
\end{theorem}

\begin{proof}
Vertical composition is composition of strategy maps. Horizontal composition sends
\[
  u:\mathcal G\Longrightarrow\mathcal G',
  \qquad
  v:\mathcal H\Longrightarrow\mathcal H'
\]
to
\[
  u\times v.
\]
The play and coplay equations follow by substitution. The response equation follows from
pseudonaturality of the stationary context decomposition, which here is equality of the explicitly
constructed morphisms.

For three games, both parenthesizations have the same uncurried play map, coplay map, and local
continuation maps after cartesian reassociation. Their classifiers are the meet of the same three
local classifiers. Associativity of meet gives the response equation for the associator.

The identity game has strategy object \(\mathbbm1\), identity play and coplay, and constantly true
classifier. Tensor is analogous, with symmetry induced by swapping strategy coordinates.
\end{proof}

\section{Open-normal strategic Mealy games}
\label{sec:open-normal-games}

\begin{definition}[Open-normal strategic Mealy game]
\label{def:open-normal-strategic-mealy-game}
An open-normal strategic Mealy game of type
\[
  (X,S)\longrightarrow(Y,R)
\]
consists of:
\begin{enumerate}
  \item a strategy object \(\Sigma\);
  \item the discrete residual category
  \[
    \operatorname{Disc}(\Sigma\times X);
  \]
  \item the totally initialized map-representable forward leg induced by
  \[
    \bar P:
    \Sigma\times X
    \longrightarrow
    (\Sigma\times X)\times Y,
    \qquad
    \bar P(\sigma,x)
    =
    ((\sigma,x),P(\sigma,x));
  \]
  \item the totally initialized map-representable backward leg induced by
  \[
    C:
    (\Sigma\times X)\times R
    \longrightarrow
    S;
  \]
  \item the extensional evaluator classified by
  \[
    \chi:
    \Sigma\times X\times R^Y\times\Sigma
    \longrightarrow
    \Omega.
  \]
\end{enumerate}
\end{definition}

A tight open-normal \(2\)-cell induced by
\[
  u:T\to\Sigma
\]
has residual map
\[
  u\times1_X:
  T\times X\to\Sigma\times X
\]
and satisfies
\[
  P_{\mathcal H}
  =
  P_{\mathcal G}\circ(u\times1_X),
\]
\[
  C_{\mathcal H}
  =
  C_{\mathcal G}
  \circ(u\times1_X\times1_R),
\]
\[
  \chi_{\mathcal H}(\tau,x,k,\tau')
  =
  \chi_{\mathcal G}(u\tau,x,k,u\tau').
\]

Define
\[
  \mathbf{OGMly}^{\mathrm{nf}}_2(\mathcal E)
\]
to have open-normal games as \(1\)-cells, tight open-normal \(2\)-cells as \(2\)-cells, and
composition given directly by the open-game play, coplay, and classifier formulas.

This is a normal-form bicategory rather than a literal composition-closed sub-bicategory of raw
optic presentations.

\subsection{Residual canonicalization}

For composable open-normal games \(\mathcal G,\mathcal H\), raw optic composition has residual
\[
  (\Sigma\times X)\times(T\times Y).
\]
The canonical composite has residual
\[
  (\Sigma\times T)\times X.
\]
Define
\[
\begin{aligned}
  h_{\mathcal H,\mathcal G}:
  (\Sigma\times T)\times X
  &\longrightarrow
  (\Sigma\times X)\times(T\times Y),
  \\
  h_{\mathcal H,\mathcal G}(\sigma,\tau,x)
  &:=
  \bigl(
    (\sigma,x),
    (\tau,P_{\mathcal G}(\sigma,x))
  \bigr).
\end{aligned}
\]

\begin{lemma}[Composite residual canonicalization]
\label{lem:composite-residual-canonicalization}
The map \(h_{\mathcal H,\mathcal G}\), together with the canonical strict initialized leg
comparisons, induces a pure residual slide from the raw optic composite to the canonical open-normal
composite. Its forward and backward maps are the ordinary composite play and coplay maps.
\end{lemma}

\begin{proof}
The raw forward residual records
\[
  (\sigma,x)
  \quad\text{and}\quad
  (\tau,P_{\mathcal G}(\sigma,x)).
\]
Precomposition by \(h_{\mathcal H,\mathcal G}\) therefore gives the canonical residual together with
the visible output
\[
  P_{\mathcal H}
  \bigl(
    \tau,
    P_{\mathcal G}(\sigma,x)
  \bigr).
\]

For the backward map, substitution gives
\[
  C_{\mathcal G}
  \Bigl(
    \sigma,
    x,
    C_{\mathcal H}
    \bigl(
      \tau,
      P_{\mathcal G}(\sigma,x),
      q
    \bigr)
  \Bigr).
\]
The total initializations are carried through the map-representable residual comparison, and the
resulting initialized leg cells are strict and invertible. Hence the resulting optic simulation is a
pure residual slide.
\end{proof}

\section{First bridge theorem: open games and open-normal Mealy games}
\label{sec:first-bridge-theorem}

Let
\[
  \mathcal G=(\Sigma,P,C,\chi)
\]
be an internal deterministic open game. Its canonical open-normal realization has residual
\(\Sigma\times X\), forward map
\[
  (\sigma,x)
  \longmapsto
  ((\sigma,x),P(\sigma,x)),
\]
backward map \(C\), and extensional evaluator classified by \(\chi\). Denote it by
\[
  J(\mathcal G).
\]

Conversely, extraction sends an open-normal strategic Mealy game to its strategy object, play map,
coplay map, and one-query classifier. Denote extraction by
\[
  U.
\]

\begin{theorem}[First bridge theorem]
\label{thm:first-bridge}
The constructions \(J\) and \(U\) define inverse symmetric monoidal isomorphisms of bicategories
\[
  \boxed{
  \mathbf{OpenGame}_2(\mathcal E)
  \cong
  \mathbf{OGMly}^{\mathrm{nf}}_2(\mathcal E).
  }
\]

Canonical residual realization also defines a locally faithful symmetric monoidal pseudofunctor
\[
  \mathcal R:
  \mathbf{OGMly}^{\mathrm{nf}}_2(\mathcal E)
  \longrightarrow
  \operatorname{StrMlyGame}^{=}_{\mathfrak C_{\mathrm{st}}}
  \bigl(
    \overline{\mathcal K}_{\mathrm{td}}
  \bigr).
\]
\end{theorem}

\begin{proof}
Extraction from the canonical realization recovers \(\Sigma,P,C,\chi\) directly.

Conversely, the defining open-normal equality
\[
  \bar P(\sigma,x)
  =
  ((\sigma,x),P(\sigma,x))
\]
shows that realization after extraction returns the same normal-form data. Thus \(J\) and \(U\) are
inverse on \(1\)-cells.

A strategy reparametrization
\[
  u:T\to\Sigma
\]
forces the residual map
\[
  u\times1_X:T\times X\to\Sigma\times X.
\]
The strict forward and backward simulation equations are precisely the play and coplay equations.
The response-strict evaluator equation is precisely the classifier equation. Hence \(J\) and \(U\)
are inverse on \(2\)-cells.

Sequential raw optic composition is identified with canonical composition by the pure residual slide
of Lemma~\ref{lem:composite-residual-canonicalization}. Since this slide is an identity after
residual coidentification, it supplies the realization compositor. The stationary decomposition and
meet evaluator give exactly the open-game response formula.

Tensor is treated by the canonical reassociation
\[
  (\Sigma\times X)\times(T\times X')
  \cong
  (\Sigma\times T)\times(X\times X').
\]
The unit, associativity, and symmetry coherence equations reduce to cartesian coherence and the
coidentifier equations.

Local faithfulness follows because residual coidentification removes pure residual-presentation
changes but does not identify distinct parameter maps or distinct strict classifier equations.
\end{proof}

\section{Branch-faithful strategic protocol implementations}
\label{sec:branch-faithful-protocols}

The optic presentation separates the forward, backward, and evaluator machines. Residual semantics
and implementation comparison are strongest when these branches are stable restrictions of one
initialized typed protocol.

\begin{definition}[Strategic protocol signature]
\label{def:strategic-protocol-signature}
A strategic protocol signature for an arena presentation \((M,\mathbf F,\mathbf G)\) and evaluator
\(\mathbf E\) consists of internal categories \(\mathsf I,\mathsf O\), fully faithful branch
functors
\[
\begin{array}{lll}
  i_F:\operatorname{dom}(F)\to\mathsf I,
  &\qquad&
  o_F:\operatorname{cod}(F)\to\mathsf O,
  \\
  i_G:\operatorname{dom}(G)\to\mathsf I,
  &&
  o_G:\operatorname{cod}(G)\to\mathsf O,
  \\
  i_E:\operatorname{dom}(E)\to\mathsf I,
  &&
  o_E:\mathsf{Truth}_\wedge\to\mathsf O,
\end{array}
\]
together with any additional protocol arrows used for phase change, observation, synchronization,
routing, and strategy-port coherence.
\end{definition}

\begin{definition}[Stable branch restriction]
\label{def:stable-branch-restriction}
Let
\[
  S:\mathsf I\rightsquigarrow\mathsf O
\]
have carrier \(P_S\). A stable branch restriction consists of a carrier subobject
\[
  m_B:P_B\hookrightarrow P_S
\]
and branch input and output subcategories such that:
\begin{enumerate}
  \item the input and output anchors of \(P_B\) factor through the branch subcategories;
  \item \(P_B\) is closed under every branch input arrow;
  \item every output produced by a branch transition lies in the branch output subcategory.
\end{enumerate}
Pullback along the branch embeddings then defines a Mealy morphism
\[
  S|_B:\mathsf I_B\rightsquigarrow\mathsf O_B.
\]
\end{definition}

\begin{definition}[Strategic protocol implementation]
\label{def:strategic-protocol-implementation}
A strategic protocol implementation of a strategic game \(\mathcal G\) consists of:
\begin{enumerate}
  \item a protocol signature
  \[
    (\mathsf I_{\mathcal G},\mathsf O_{\mathcal G});
  \]
  \item a transition-closed initialized Mealy morphism
  \[
    \mathsf{SProt}(\mathcal G):
    \mathsf I_{\mathcal G}
    \rightsquigarrow
    \mathsf O_{\mathcal G};
  \]
  \item stable carrier subobjects
  \[
    P_F\hookrightarrow P_{\mathsf{SProt}},
    \qquad
    P_G\hookrightarrow P_{\mathsf{SProt}},
    \qquad
    P_E\hookrightarrow P_{\mathsf{SProt}};
  \]
  \item invertible initialized Mealy simulations identifying the branch restrictions with the
  forward leg, backward leg, and response evaluator;
  \item a distinguished ready-state family
  \[
    \rho_{\mathcal G}:R_{\mathcal G}\to J_{\mathsf{SProt}}.
  \]
\end{enumerate}
\end{definition}

The stable carrier subobjects are essential. Shared phase objects do not by themselves determine to
which branch a machine state belongs.

\section{Structured canonical protocols for open-normal games}
\label{sec:canonical-protocols}

Let
\[
  \mathcal G=(\Sigma,P,C,\chi):(X,S)\to(Y,R)
\]
be an open-normal strategic Mealy game. Put
\[
  \mathsf C_{\mathcal G}:=X\times R^Y.
\]

\subsection{Existence of the protocol categories}

A Grothendieck topos has free internal categories on internal graphs. Explicitly, if
\[
  G_1
  \mathrel{\substack{\xrightarrow{\ s\ }\\[-1mm]\xrightarrow[\ t\ ]{}}}
  G_0
\]
is an internal graph, the object of length-\(n\) composable paths is an iterated pullback
\(G_1^{[n]}\), and
\[
  \coprod_{n\in\mathbb N}G_1^{[n]}
\]
carries the free internal category structure by concatenation. Presentations with specified internal
subcategory inclusions are obtained by the chosen pushouts and effective quotients in
\(\operatorname{Cat}(\mathcal E)\).

\subsection{Protocol input category and strategy ports}

Begin with the coproduct of internal categories
\[
\begin{aligned}
  &
  \operatorname{Disc}(\Sigma)
  +
  \operatorname{Disc}(\Sigma\times X)
  +
  \operatorname{Disc}(\Sigma\times X\times R)
  \\
  &\hspace{35mm}
  +
  \mathsf{Qry}
  \bigl(
    \Sigma,
    \operatorname{Disc}(\mathsf C_{\mathcal G})
  \bigr).
\end{aligned}
\]
Write the tagged object families as
\[
  \mathsf p_\sigma,
  \qquad
  \mathsf f_{\sigma,x},
  \qquad
  \mathsf g_{\sigma,x,r},
  \qquad
  \mathsf e_{\sigma,c}.
\]

Freely adjoin a return arrow
\[
  \rho_{\sigma,x,r}:
  \mathsf g_{\sigma,x,r}
  \longrightarrow
  \mathsf f_{\sigma,x}
\]
for every \((\sigma,x,r)\).

Freely adjoin strategy-tether arrows
\[
  \tau^F_{\sigma,x}:
  \mathsf f_{\sigma,x}
  \longrightarrow
  \mathsf p_\sigma,
\]
\[
  \tau^G_{\sigma,x,r}:
  \mathsf g_{\sigma,x,r}
  \longrightarrow
  \mathsf p_\sigma,
\]
and
\[
  \tau^E_{\sigma,c}:
  \mathsf e_{\sigma,c}
  \longrightarrow
  \mathsf p_\sigma.
\]

The resulting internal category is
\[
  \mathsf I_{\mathcal G}.
\]

There is a strategy projection functor
\[
  \pi_{\mathcal G}:
  \mathsf I_{\mathcal G}
  \longrightarrow
  \operatorname{Disc}(\Sigma)
\]
sending every tagged object with strategy coordinate \(\sigma\) to \(\sigma\), and sending every
return, tether, and query arrow to the identity at its incumbent strategy.

Let
\[
  \mathsf{OneQry}_{\mathcal G}
  :=
  \Sigma\times\mathsf C_{\mathcal G}\times\Sigma
\]
be the subobject of one-letter evaluator query arrows. It carries the named projections
\[
  \operatorname{inc}_{\mathcal G},
  \operatorname{alt}_{\mathcal G}:
  \mathsf{OneQry}_{\mathcal G}
  \rightrightarrows
  \Sigma
\]
to the incumbent and alternative strategy coordinates.

Under target-to-source processing, a state anchored at
\(\mathsf f_{\sigma,x}\) processes \(\rho_{\sigma,x,r}\) and moves to a backward state anchored at
\(\mathsf g_{\sigma,x,r}\).

The tether arrows have target \(\mathsf p_\sigma\). No protocol carrier state is anchored at a
strategy-port object, so tether arrows and all paths with target in the strategy-port branch have
empty transition fibres. Their purpose is to constrain protocol reparametrizations without adding
execution behaviour.

\subsection{Protocol output category}

Begin with
\[
  \operatorname{Disc}((\Sigma\times X)\times Y)
  +
  \operatorname{Disc}(S)
  +
  \mathsf{Truth}_\wedge.
\]
Write its tagged object families as
\[
  \mathsf F_{((\sigma,x),y)},
  \qquad
  \mathsf G_s,
  \qquad
  \mathsf E_*.
\]
Freely adjoin arrows
\[
  \bar\rho_{\sigma,x,r}:
  \mathsf G_{C(\sigma,x,r)}
  \longrightarrow
  \mathsf F_{((\sigma,x),P(\sigma,x))}.
\]
The resulting internal category is
\[
  \mathsf O_{\mathcal G}.
\]

\begin{proposition}[Existence of structured canonical protocol signatures]
\label{prop:canonical-protocol-signature-exists}
The input and output presentations determine internal categories
\[
  \mathsf I_{\mathcal G},
  \qquad
  \mathsf O_{\mathcal G}.
\]
Their stated branch inclusions are fully faithful. The strategy projection and one-query projections
are internal morphisms with the stated types.
\end{proposition}

\begin{proof}
Apply the chosen free-internal-category functor to the displayed internal graphs and impose the
existing composition equations of the evaluator branch using the chosen effective quotient.

Disjointness of coproducts makes each branch inclusion monic on objects and arrows. The new return
and tether generators create no additional arrows between two objects of the same original branch.
Indeed, no generator leaves a strategy-port object, and the only nonidentity generator between the
forward and backward families is the return arrow directed from the backward tag to the forward tag.
Hence no new path has both endpoints in one original branch unless it already belongs to that
branch. The branch inclusions are therefore fully faithful.

The strategy projection and the one-query coordinate maps are defined on generators and preserve
all imposed equations, so they descend to the chosen quotient.
\end{proof}

\subsection{Protocol carrier and dynamics}

Define the carrier
\[
  Z_{\mathcal G}
  :=
  (\Sigma\times X)
  +
  (\Sigma\times X\times R)
  +
  (\Sigma\times\mathsf C_{\mathcal G}).
\]
The three coproduct summands are the forward, backward, and evaluator branches.

The input anchors are
\[
  (\sigma,x)\longmapsto\mathsf f_{\sigma,x},
\]
\[
  (\sigma,x,r)\longmapsto\mathsf g_{\sigma,x,r},
\]
\[
  (\sigma,c)\longmapsto\mathsf e_{\sigma,c}.
\]
There is no carrier state over any \(\mathsf p_\sigma\).

The output anchors are
\[
  (\sigma,x)
  \longmapsto
  \mathsf F_{((\sigma,x),P(\sigma,x))},
\]
\[
  (\sigma,x,r)
  \longmapsto
  \mathsf G_{C(\sigma,x,r)},
\]
\[
  (\sigma,c)
  \longmapsto
  \mathsf E_*.
\]

On a return generator, define
\[
  (\sigma,x)
  \xrightarrow{
    \rho_{\sigma,x,r}/\bar\rho_{\sigma,x,r}
  }
  (\sigma,x,r).
\]
On an evaluator query word
\[
  \sigma'_1\cdots\sigma'_n
\]
at context \(c\), the evaluator state remains \((\sigma,c)\) and emits
\[
  \bigwedge_{i=1}^{n}
  \chi(\sigma,c,\sigma'_i).
\]
Identity input arrows act identically.

Every remaining nonidentity path with nonempty transition fibre is a path in the evaluator branch
or a return arrow composed with identities. The preceding assignments therefore extend uniquely to
all admissible paths.

\begin{definition}[Structured canonical branch-faithful protocol]
\label{def:canonical-branch-faithful-protocol}
The resulting totally initialized Mealy morphism, together with its named strategy-port and
one-query structure, is
\[
  \mathsf{Prot}(\mathcal G):
  \mathsf I_{\mathcal G}
  \rightsquigarrow
  \mathsf O_{\mathcal G}.
\]
Its distinguished ready family is the coproduct inclusion
\[
  R_{\mathcal G}
  :=
  (\Sigma\times X)
  +
  (\Sigma\times\mathsf C_{\mathcal G})
  \longrightarrow
  Z_{\mathcal G}.
\]
Backward states are not independently ready; they are reached as derivatives of forward ready states
along return arrows.
\end{definition}

\begin{proposition}[Canonical branch identification]
\label{prop:canonical-branch-identification}
The three carrier coproduct summands are stable branch restrictions. Their restrictions are:
\begin{enumerate}
  \item the canonical forward map
  \[
    (\sigma,x)
    \longmapsto
    ((\sigma,x),P(\sigma,x));
  \]
  \item the canonical backward map
  \[
    (\sigma,x,r)
    \longmapsto
    C(\sigma,x,r);
  \]
  \item the extensional response evaluator classified by \(\chi\).
\end{enumerate}
\end{proposition}

\begin{proof}
The return arrows are not arrows of any branch subcategory. Tether arrows have no carrier state over
their targets. Hence each carrier summand is closed under the arrows of its designated branch
category.

The restrictions of the anchors, transitions, and outputs are exactly the three displayed
constructions.
\end{proof}

\subsection{Structured canonical protocol simulations}

Let
\[
  \mathcal G=(\Sigma,P_{\mathcal G},C_{\mathcal G},\chi_{\mathcal G}),
\]
\[
  \mathcal H=(T,P_{\mathcal H},C_{\mathcal H},\chi_{\mathcal H})
\]
have the same signed boundary.

\begin{definition}[Structured protocol functor]
\label{def:structured-protocol-functor}
A structured protocol functor from the signature of \(\mathcal H\) to the signature of
\(\mathcal G\) consists of internal functors
\[
  I:\mathsf I_{\mathcal H}\to\mathsf I_{\mathcal G},
  \qquad
  O:\mathsf O_{\mathcal H}\to\mathsf O_{\mathcal G}
\]
such that:
\begin{enumerate}
  \item \(I\) preserves the parameter-port, forward, backward, and evaluator object families;
  \item \(I\) preserves every non-strategy coordinate;
  \item the restriction of \(I\) to the parameter-port branch is
  \[
    \operatorname{Disc}(u):
    \operatorname{Disc}(T)\to\operatorname{Disc}(\Sigma)
  \]
  for a map \(u:T\to\Sigma\);
  \item \(I\) preserves return and tether generator families;
  \item on one-letter query generators, \(I\) preserves both incumbent and alternative projections:
  \[
    \operatorname{inc}_{\mathcal G}I
    =
    u\,\operatorname{inc}_{\mathcal H},
    \qquad
    \operatorname{alt}_{\mathcal G}I
    =
    u\,\operatorname{alt}_{\mathcal H};
  \]
  \item \(O\) preserves the forward, backward, and truth output branches;
  \item \(O\) preserves all non-strategy output coordinates and every return generator.
\end{enumerate}
\end{definition}

For every map
\[
  u:T\to\Sigma,
\]
define
\[
  I_u:
  \mathsf I_{\mathcal H}
  \longrightarrow
  \mathsf I_{\mathcal G}
\]
by
\[
  \mathsf p_\tau\longmapsto\mathsf p_{u\tau},
\]
\[
  \mathsf f_{\tau,x}\longmapsto\mathsf f_{u\tau,x},
\]
\[
  \mathsf g_{\tau,x,r}\longmapsto\mathsf g_{u\tau,x,r},
\]
\[
  \mathsf e_{\tau,c}\longmapsto\mathsf e_{u\tau,c},
\]
and by applying \(u\) to every strategy letter of a query word, every return generator, and every
tether generator.

The input functor \(I_u\) exists for every \(u\) by the freeness of the input presentation.

Define a proposed output object map by
\[
  \mathsf F_{((\tau,x),y)}
  \longmapsto
  \mathsf F_{((u\tau,x),y)},
\]
\[
  \mathsf G_s\longmapsto\mathsf G_s,
  \qquad
  \mathsf E_*\longmapsto\mathsf E_*.
\]

\begin{lemma}[Existence of the output reindexing functor]
\label{lem:output-reindexing-exists}
The preceding object assignment extends uniquely to a branch- and return-preserving internal functor
\[
  O_u:
  \mathsf O_{\mathcal H}
  \longrightarrow
  \mathsf O_{\mathcal G}
\]
if and only if
\[
  P_{\mathcal H}(\tau,x)
  =
  P_{\mathcal G}(u\tau,x)
\]
and
\[
  C_{\mathcal H}(\tau,x,r)
  =
  C_{\mathcal G}(u\tau,x,r).
\]
When these equations hold,
\[
  O_u
  \bigl(
    \bar\rho^{\mathcal H}_{\tau,x,r}
  \bigr)
  :=
  \bar\rho^{\mathcal G}_{u\tau,x,r}.
\]
\end{lemma}

\begin{proof}
The source return generator is
\[
  \bar\rho^{\mathcal H}_{\tau,x,r}:
  \mathsf G_{C_{\mathcal H}(\tau,x,r)}
  \longrightarrow
  \mathsf F_{((\tau,x),P_{\mathcal H}(\tau,x))}.
\]
A return-preserving output functor must send it to
\[
  \bar\rho^{\mathcal G}_{u\tau,x,r}:
  \mathsf G_{C_{\mathcal G}(u\tau,x,r)}
  \longrightarrow
  \mathsf F_{((u\tau,x),P_{\mathcal G}(u\tau,x))}.
\]
The source and target of this arrow agree with the prescribed images of the source and target of the
original generator precisely when the two displayed equations hold.

When they hold, the assignment is defined on every generator and preserves all relations, so it
extends uniquely to the freely presented output category.
\end{proof}

Assume from now on that the play and coplay equations hold, so \(O_u\) exists.

Input reindexing and output postcomposition give Mealy morphisms with common source and target
interfaces:
\[
  I_u^*\mathsf{Prot}(\mathcal G),
  \qquad
  (O_u)_\#\circ\mathsf{Prot}(\mathcal H).
\]

Put
\[
\begin{aligned}
  \bar u:
  Z_{\mathcal H}
  &\longrightarrow
  Z_{\mathcal G},
  \\
  \bar u
  &:=
  (u\times1_X)
  +
  (u\times1_X\times1_R)
  +
  (u\times1_{\mathsf C_{\mathcal G}}).
\end{aligned}
\]

The carrier of \(I_u^*\mathsf{Prot}(\mathcal G)\) is
\[
  (\mathsf I_{\mathcal H})_0
  \times_{I_{u,0},(\mathsf I_{\mathcal G})_0,p_{\mathcal G}}
  Z_{\mathcal G}.
\]

Define
\[
\begin{aligned}
  \kappa_u:
  Z_{\mathcal H}
  &\longrightarrow
  (\mathsf I_{\mathcal H})_0
  \times_{(\mathsf I_{\mathcal G})_0}
  Z_{\mathcal G},
  \\
  z
  &\longmapsto
  \bigl(
    p_{\mathcal H}(z),
    \bar u(z)
  \bigr).
\end{aligned}
\]

\begin{definition}[Structured canonical protocol simulation]
\label{def:canonical-protocol-simulation}
A structured canonical protocol simulation induced by \(u\) is the strict initialized Mealy
simulation
\[
  I_u^*\mathsf{Prot}(\mathcal G)
  \Longrightarrow
  (O_u)_\#\circ\mathsf{Prot}(\mathcal H)
\]
whose carrier comparison is \(\kappa_u\).
\end{definition}

\begin{proposition}[Canonical protocol simulation equations]
\label{prop:canonical-protocol-simulation-equations}
Let \(u:T\to\Sigma\).

The output functor \(O_u\) exists if and only if
\[
  P_{\mathcal H}(\tau,x)
  =
  P_{\mathcal G}(u\tau,x)
\]
and
\[
  C_{\mathcal H}(\tau,x,r)
  =
  C_{\mathcal G}(u\tau,x,r).
\]

After \(O_u\) has been constructed, the map \(\kappa_u\) defines a strict initialized protocol
simulation if and only if
\[
  \chi_{\mathcal H}(\tau,c,\tau')
  =
  \chi_{\mathcal G}(u\tau,c,u\tau').
\]

Conversely, every structured protocol functor together with a branch-preserving strict initialized
simulation between structured canonical protocols is induced by a unique map \(u:T\to\Sigma\).
\end{proposition}

\begin{proof}
The existence statement for \(O_u\) is
Lemma~\ref{lem:output-reindexing-exists}.

The pullback equation
\[
  I_{u,0}p_{\mathcal H}(z)
  =
  p_{\mathcal G}\bar u(z)
\]
holds on each carrier summand, so \(\kappa_u\) is defined.

On a return generator,
\[
  \bar u
  \delta_{\mathcal H}
  \bigl(
    \rho^{\mathcal H}_{\tau,x,r},
    (\tau,x)
  \bigr)
  =
  \bar u(\tau,x,r)
  =
  (u\tau,x,r),
\]
while
\[
  \delta_{\mathcal G}
  \bigl(
    I_u(\rho^{\mathcal H}_{\tau,x,r}),
    \bar u(\tau,x)
  \bigr)
  =
  \delta_{\mathcal G}
  \bigl(
    \rho^{\mathcal G}_{u\tau,x,r},
    (u\tau,x)
  \bigr)
  =
  (u\tau,x,r).
\]
Thus the transition equation holds on return generators. It is immediate on identities and
evaluator query arrows. Tether arrows have empty transition fibres.

The output on a return generator is preserved because
\[
  O_u
  \bigl(
    \bar\rho^{\mathcal H}_{\tau,x,r}
  \bigr)
  =
  \bar\rho^{\mathcal G}_{u\tau,x,r}.
\]

On a one-letter evaluator query, strict output preservation is exactly
\[
  \chi_{\mathcal H}(\tau,c,\tau')
  =
  \chi_{\mathcal G}(u\tau,c,u\tau').
\]
For a finite query word, equality follows by taking finite meets of the one-letter equations.
Ready-state preservation is immediate.

Conversely, restriction of the structured input functor to the strategy-port branch determines a
unique map
\[
  u:T\to\Sigma.
\]
Preservation of the tether
\[
  \mathsf f_{\tau,x}\to\mathsf p_\tau
\]
forces
\[
  \mathsf f_{\tau,x}\longmapsto\mathsf f_{u\tau,x},
\]
because the non-strategy coordinate \(x\) is preserved and the only forward tether with target
\(\mathsf p_{u\tau}\) is the displayed one. The backward and evaluator tethers similarly force
\[
  \mathsf g_{\tau,x,r}\longmapsto\mathsf g_{u\tau,x,r},
\]
\[
  \mathsf e_{\tau,c}\longmapsto\mathsf e_{u\tau,c}.
\]
Preservation of the incumbent and alternative one-query projections forces
\[
  ([\tau'],1_c)
  \longmapsto
  ([u\tau'],1_c).
\]
Preservation of return generators then forces
\[
  \rho_{\tau,x,r}
  \longmapsto
  \rho_{u\tau,x,r}.
\]
By freeness, the input functor is \(I_u\).

Each canonical branch has exactly one carrier state over each carrier tag. Hence the carrier
comparison is necessarily \(\kappa_u\). Preservation of output return generators gives the play and
coplay equations, and strictness on one-letter evaluator queries gives the classifier equation.
\end{proof}

\subsection{Direct composition of structured canonical protocols}

Let
\[
  \mathsf P:
  (X,S)\to(Y,R),
  \qquad
  \mathsf Q:
  (Y,R)\to(Z,Q)
\]
be structured canonical protocols. Extract their unique data
\[
  (\Sigma,P_{\mathsf P},C_{\mathsf P},\chi_{\mathsf P})
\]
and
\[
  (T,P_{\mathsf Q},C_{\mathsf Q},\chi_{\mathsf Q}).
\]

Define
\[
  \mathsf Q\odot\mathsf P
\]
to be the structured canonical protocol compiled from the strategy object \(\Sigma\times T\) and
the maps
\[
  P_{\mathsf Q\odot\mathsf P}
  ((\sigma,\tau),x)
  =
  P_{\mathsf Q}
  \bigl(
    \tau,
    P_{\mathsf P}(\sigma,x)
  \bigr),
\]
\[
\begin{aligned}
  C_{\mathsf Q\odot\mathsf P}
  ((\sigma,\tau),x,q)
  =
  C_{\mathsf P}
  \Bigl(
    \sigma,
    x,
    C_{\mathsf Q}
    \bigl(
      \tau,
      P_{\mathsf P}(\sigma,x),
      q
    \bigr)
  \Bigr),
\end{aligned}
\]
and
\[
\begin{aligned}
&\chi_{\mathsf Q\odot\mathsf P}
\bigl(
  (\sigma,\tau),
  x,
  k,
  (\sigma',\tau')
\bigr)
\\
&\qquad=
\chi_{\mathsf P}
\bigl(
  \sigma,
  x,
  \bar\kappa_{\mathsf Q}(\tau,k),
  \sigma'
\bigr)
\\
&\qquad\qquad\wedge
\chi_{\mathsf Q}
\bigl(
  \tau,
  P_{\mathsf P}(\sigma,x),
  k,
  \tau'
\bigr).
\end{aligned}
\]

Identity protocols and tensor protocols are defined by compiling the corresponding direct formulas.

The external source and target of a structured canonical protocol are its signed boundaries. Its
internal protocol categories are not the objects used for bicategorical composition. The operation
\(\odot\) is compiler composition and is not ordinary Mealy pipeline composition of the compiled
machines.

\begin{theorem}[Structured canonical protocol bicategory]
\label{thm:canonical-protocol-bicategory}
Signed pairs of objects, structured canonical protocols, and structured canonical protocol
simulations form a symmetric monoidal bicategory
\[
  \mathbf{CanProt}^{\mathrm{str}}_2(\mathcal E)
\]
under the directly defined composition \(\odot\).
\end{theorem}

\begin{proof}
The extracted play maps of the two parenthesizations of a triple composite are equal after cartesian
reassociation. The same is true of the coplay maps by associativity of substitution.

The three local stationary contexts in either parenthesization are equal by
Proposition~\ref{prop:stationary-context-doctrine}. Hence both classifiers are the meet of the same
three local classifier values. Associativity of meet identifies them.

The canonical strategy reassociation map therefore defines a structured strict protocol simulation
between the two compiled parenthesizations. Its output functor exists because the play and coplay
maps agree after reassociation. Its classifier equation follows from associativity of meet. Its
inverse is induced by the opposite reassociation. The pentagon follows from cartesian coherence.

Identity and tensor are analogous. Symmetry is induced by the cartesian swap of strategy objects and
the corresponding swaps of protocol tags and ports. Horizontal composition of protocol simulations
is product of their strategy maps. Interchange follows from functoriality of product.
\end{proof}

\section{Second bridge theorem: open-normal games and structured canonical protocols}
\label{sec:second-bridge-theorem}

Let
\[
  \mathsf{Prot}:
  \mathbf{OGMly}^{\mathrm{nf}}_2(\mathcal E)
  \longrightarrow
  \mathbf{CanProt}^{\mathrm{str}}_2(\mathcal E)
\]
be the chosen structured canonical protocol compiler.

Let
\[
  \mathsf{Ext}:
  \mathbf{CanProt}^{\mathrm{str}}_2(\mathcal E)
  \longrightarrow
  \mathbf{OGMly}^{\mathrm{nf}}_2(\mathcal E)
\]
extract the strategy object and the unique forward, backward, and evaluator branch maps.

\begin{theorem}[Second bridge theorem]
\label{thm:second-bridge}
Structured canonical protocol construction and branch extraction are inverse symmetric monoidal
isomorphisms of bicategories:
\[
  \boxed{
  \mathbf{OGMly}^{\mathrm{nf}}_2(\mathcal E)
  \cong
  \mathbf{CanProt}^{\mathrm{str}}_2(\mathcal E).
  }
\]
Consequently,
\[
  \boxed{
  \mathbf{OpenGame}_2(\mathcal E)
  \cong
  \mathbf{CanProt}^{\mathrm{str}}_2(\mathcal E).
  }
\]
\end{theorem}

\begin{proof}
By Proposition~\ref{prop:canonical-branch-identification}, extraction from
\(\mathsf{Prot}(\mathcal G)\) recovers the original strategy object, play map, coplay map, and
classifier.

Conversely, a structured canonical protocol includes the chosen compiler presentation and is
determined by its extracted tuple together with the named branch, port, generator, and ready-family
structure. Recompiling from the extracted maps therefore returns the same structured protocol
literally.

Proposition~\ref{prop:canonical-protocol-simulation-equations} gives the corresponding inverse
statement on \(2\)-cells. In particular, the output functor of a structured protocol simulation is
constructed only after the play and coplay equations have been established.

The directly defined protocol composition uses exactly the open-normal play, coplay, and response
formulas. Hence
\[
  \mathsf{Prot}(\mathcal H\circ\mathcal G)
  =
  \mathsf{Prot}(\mathcal H)
  \odot
  \mathsf{Prot}(\mathcal G)
\]
up to the canonical cartesian reassociations built into the bicategories. Identity, tensor, and
symmetry are preserved for the same reason.
\end{proof}

\subsection{The replete structured protocol class}

For many purposes the literal chosen presentation of a canonical protocol is inessential, but the
branch, port, generator, and ready-family markings cannot all be discarded: they are required to
define extraction and compiler composition.

\begin{definition}[Replete structured canonical protocol]
\label{def:replete-structured-canonical-protocol}
A replete structured canonical protocol is a strategic protocol implementation carrying:
\begin{enumerate}
  \item named forward, backward, and evaluator branch embeddings;
  \item stable branch carrier subobjects;
  \item a strategy-port projection;
  \item named return and tether generator families;
  \item named one-query projections;
  \item a distinguished ready family;
\end{enumerate}
and which is structurally isomorphic to a chosen canonical protocol.

A morphism of replete structured protocols is a structure-preserving initialized protocol
simulation.
\end{definition}

Let
\[
  \mathbf{CanProt}^{\mathrm{rep}}_2(\mathcal E)
\]
be the bicategory whose composition is defined by extraction, open-normal composition, and
recompilation.

\begin{proposition}[Replete presentation biequivalence]
\label{prop:replete-presentation-biequivalence}
The inclusion
\[
  \mathbf{CanProt}^{\mathrm{str}}_2(\mathcal E)
  \longrightarrow
  \mathbf{CanProt}^{\mathrm{rep}}_2(\mathcal E)
\]
is a symmetric monoidal biequivalence.
\end{proposition}

\begin{proof}
Every replete structured protocol is, by definition, structurally isomorphic to a chosen canonical
protocol. Thus the inclusion is essentially surjective on each hom-category.

A structure-preserving simulation between chosen canonical protocols is uniquely determined by its
strategy map by Proposition~\ref{prop:canonical-protocol-simulation-equations}. The same statement
holds after transport along structural isomorphisms. Hence the inclusion is locally fully faithful
up to equivalence.

Composition and tensor in the replete bicategory are defined through extraction and recompilation,
so the inclusion preserves them up to the canonical structural isomorphisms.
\end{proof}

The proposition does not discard the structural markings. If all branch, port, generator, and
ready-family data are forgotten, branch extraction and compiler composition are no longer defined
from the underlying Mealy machines alone.

\section{Strategic residual semantics and minimization}
\label{sec:strategic-residual-semantics}

Let
\[
  \mathsf S=
  \mathsf{SProt}(\mathcal G):
  \mathsf I\rightsquigarrow\mathsf O
\]
be a strategic protocol implementation with carrier \(P\) and distinguished ready map
\[
  \rho:R\longrightarrow J_{\mathsf S}.
\]
Put
\[
  i_{\mathcal G}
  :=
  j_{\mathsf S}\rho:
  R\longrightarrow P.
\]
Chapter~3 gives the behaviour assignment
\[
  \beta_P:
  P
  \longrightarrow
  \operatorname{Beh}_0(\mathsf I,\mathsf O).
\]

\begin{definition}[Strategic behaviour specification]
\label{def:strategic-behaviour-specification}
The strategic behaviour specification is
\[
  k_{\mathcal G}
  :=
  \beta_Pi_{\mathcal G}:
  R
  \longrightarrow
  \operatorname{Beh}_0(\mathsf I,\mathsf O).
\]
\end{definition}

Define the ready anchors
\[
  p_R:=p_Pi_{\mathcal G}:R\to\mathsf I_0,
  \qquad
  q_R:=q_Pi_{\mathcal G}:R\to\mathsf O_0.
\]
Its formal derivative object is
\[
  D_R
  :=
  \mathsf I_1
  \times_{t_{\mathsf I},\mathsf I_0,p_R}
  R.
\]
Residual evaluation is
\[
  d_R(a,x)
  :=
  \partial_a k_{\mathcal G}(x).
\]

\begin{definition}[Strategic Nerode relation]
\label{def:strategic-nerode-relation}
The strategic Nerode relation
\[
  {\equiv_{\mathcal G}}
  \rightrightarrows
  D_R
\]
is the kernel pair of \(d_R\) in the slice over
\[
  \mathsf I_0\times\mathsf O_0.
\]
\end{definition}

Because the three branches belong to one typed protocol, the relation simultaneously observes
forward execution, continuation-return behaviour, and counterfactual response.

\begin{theorem}[Minimal strategic realization]
\label{thm:minimal-strategic-realization}
The following hold:
\begin{enumerate}
  \item the effective quotient
  \[
    D_R/{\equiv_{\mathcal G}}
  \]
  exists;
  \item it is canonically isomorphic to the derivative image
  \[
    \operatorname{Der}_0(k_{\mathcal G});
  \]
  \item the restriction of the canonical behaviour machine to this image is the minimal residual
  realization
  \[
    \operatorname{Min}(k_{\mathcal G});
  \]
  \item the reachable part of \(\mathsf S\) maps by a regular-epimorphic strict semantic
  homomorphism onto \(\operatorname{Min}(k_{\mathcal G})\);
  \item every reachable behaviourally separated realization of \(k_{\mathcal G}\) is isomorphic to
  \(\operatorname{Min}(k_{\mathcal G})\).
\end{enumerate}
\end{theorem}

\begin{proof}
A Grothendieck topos is Barr-exact and locally cartesian closed, so Assumption~3.1.1 holds. Apply the
categorical Myhill--Nerode theorem of Chapter~3 to
\[
  k_{\mathcal G}
  =
  \beta_Pi_{\mathcal G}.
\]
That theorem identifies the effective Nerode quotient with the derivative image, equips it with the
restricted canonical Mealy structure, supplies the regular-epimorphic semantic quotient from the
reachable part, and proves uniqueness among reachable behaviourally separated realizations.
\end{proof}

\subsection{One-query and full evaluator observations}

Let
\[
  E:
  \mathsf{Qry}(\Sigma,\mathcal C)
  \rightsquigarrow
  \mathsf{Truth}_\wedge
\]
be the evaluator branch of a strategic protocol.

For an evaluator state \(z\) over \((\sigma,c)\), define its one-query profile by
\[
  \operatorname{obs}_1(z)(\sigma')
  :=
  \omega_E
  \bigl(
    ([\sigma'],1_c),
    z
  \bigr).
\]
Internally this gives a morphism
\[
  \operatorname{obs}_1:
  P_E
  \longrightarrow
  \operatorname{Prof}_{\Sigma,\mathcal C},
\]
where the fibre of
\(\operatorname{Prof}_{\Sigma,\mathcal C}\) over \((\sigma,c)\) is \(\Omega^\Sigma\).

Evaluation of complete evaluator behaviour on one-letter deviations defines
\[
  \operatorname{ev}_1:
  \operatorname{Beh}_0
  \bigl(
    \mathsf{Qry}(\Sigma,\mathcal C),
    \mathsf{Truth}_\wedge
  \bigr)
  \longrightarrow
  \operatorname{Prof}_{\Sigma,\mathcal C}
\]
such that
\[
  \operatorname{obs}_1
  =
  \operatorname{ev}_1\beta_E.
\]

Let
\[
  D_E
\]
be the formal derivative object generated from the evaluator ready family. Define
\[
  \equiv^E_{\mathrm{full}}
  :=
  \ker(d_E),
\]
and
\[
  \equiv^E_1
  :=
  \ker(\operatorname{ev}_1d_E).
\]

\begin{proposition}[One-query comparison]
\label{prop:one-query-comparison}
There is an inclusion of internal equivalence relations
\[
  \equiv^E_{\mathrm{full}}
  \subseteq
  \equiv^E_1.
\]
If \(E\) is the canonical extensional evaluator, then
\[
  \equiv^E_{\mathrm{full}}
  =
  \equiv^E_1.
\]
\end{proposition}

\begin{proof}
The inclusion follows from the factorization
\[
  \operatorname{ev}_1d_E.
\]

For the extensional evaluator, let
\[
  \operatorname{Prof}^{\mathrm{inv}}_{\Sigma,\mathcal C}
  \hookrightarrow
  \operatorname{Prof}_{\Sigma,\mathcal C}
\]
be the subobject of context-isomorphism-invariant profiles. Define
\[
  \operatorname{Ext}:
  \operatorname{Prof}^{\mathrm{inv}}_{\Sigma,\mathcal C}
  \longrightarrow
  \operatorname{Beh}_0
  \bigl(
    \mathsf{Qry}(\Sigma,\mathcal C),
    \mathsf{Truth}_\wedge
  \bigr)
\]
by
\[
  \operatorname{Ext}(\varphi)(\sigma'_1\cdots\sigma'_n,f)
  :=
  \bigwedge_{i=1}^{n}\varphi(\sigma'_i).
\]
Context invariance makes this independent of transport along \(f\).

Then
\[
  \operatorname{ev}_1\operatorname{Ext}
  =
  1.
\]
On the extensional behaviour image,
\[
  \operatorname{Ext}\operatorname{ev}_1
  =
  1.
\]
Hence \(\operatorname{ev}_1\) is monic on extensional evaluator behaviours, and the two kernel pairs
coincide.
\end{proof}

For the complete branch-faithful protocol, full strategic behaviour also includes forward and
backward branch observations. Proposition~\ref{prop:one-query-comparison} therefore concerns the
evaluator derivative subobject. Complete protocol minimization restricts on evaluator derivatives to
a refinement of one-query response minimization, and the two coincide for canonical extensional
evaluator derivatives.

\subsection{Reachable closure from a ready family}

The following construction will also be used for history-indexed residual games.

\begin{lemma}[Reachable closure generated by selected states]
\label{lem:reachable-closure-selected-states}
Let
\[
  P:A\rightsquigarrow B
\]
be a Mealy morphism and let
\[
  i:K\to P_0
\]
be a family of selected states.

Define the orbit object
\[
  \mathcal O_{P,i}
  :=
  A_1\times_{t_A,A_0,p_Pi}K
\]
and orbit map
\[
  \operatorname{orb}_{P,i}(a,k)
  :=
  \delta_P(a,i(k)).
\]
Factor it as
\[
  \mathcal O_{P,i}
  \twoheadrightarrow
  \operatorname{Reach}_0(P,i)
  \xhookrightarrow{m_i}
  P_0.
\]

Then \(\operatorname{Reach}_0(P,i)\) carries a unique Mealy submorphism
\[
  \operatorname{Reach}(P,i):A\rightsquigarrow B
\]
for which \(m_i\) is a strict Mealy homomorphism. The selected map \(i\) factors through
\(\operatorname{Reach}_0(P,i)\).
\end{lemma}

\begin{proof}
A representative reachable state has the form
\[
  \delta_P(a,i(k)).
\]
If \(b\) is an input arrow composable at this state, then
\[
  \delta_P
  \bigl(
    b,\delta_P(a,i(k))
  \bigr)
  =
  \delta_P(a\circ b,i(k)),
\]
which again lies in the image of the orbit map.

Categorically, pull back the regular epimorphism
\[
  \mathcal O_{P,i}
  \twoheadrightarrow
  \operatorname{Reach}_0(P,i)
\]
along the domain of the proposed restricted transition. On that regular-epimorphic cover, the
transition factors through \(m_i\) by the multiplication equation above. Regular-epimorphic descent
and monicity of \(m_i\) give a unique map
\[
  \delta_{\operatorname{Reach}}:
  A_1\times_{t_A,A_0,p_{\operatorname{Reach}}}
  \operatorname{Reach}_0(P,i)
  \to
  \operatorname{Reach}_0(P,i)
\]
such that
\[
  m_i\delta_{\operatorname{Reach}}
  =
  \delta_P(1\times m_i).
\]

Define
\[
  \omega_{\operatorname{Reach}}
  :=
  \omega_P(1\times m_i).
\]
The Mealy unit and multiplication laws follow after composition with the monomorphism \(m_i\).

Finally,
\[
  i(k)
  =
  \delta_P(1_{p_Pi(k)},i(k)),
\]
so \(i\) belongs to the image of the orbit map and factors through \(m_i\).
\end{proof}

\section{History-indexed residual games and subgame perfection}
\label{sec:history-indexed-residual-games}

A branch-faithful protocol has distinct carrier branches. A single state need not simultaneously
encode the forward, backward, and evaluator residual states. Game-theoretic residuals therefore
require explicit synchronization over both a strategy profile and a legal history.

\begin{definition}[History-indexed strategic implementation]
\label{def:history-indexed-strategic-implementation}
A history-indexed strategic implementation with strategy object \(\Sigma\) consists of:
\begin{enumerate}
  \item an internal object \(H\) of legal histories and a root
  \[
    h_0:\mathbbm1\to H;
  \]
  \item an internal object \(L\) of legal one-step extensions with origin and destination maps
  \[
    o_L,d_L:L\rightrightarrows H;
  \]
  \item ready-state maps
  \[
    \rho_F:\Sigma\times H\to J_F,
    \qquad
    \rho_G:\Sigma\times H\to J_G,
    \qquad
    \rho_E:\Sigma\times H\to J_E;
  \]
  \item for each branch \(B\in\{F,G,E\}\), an input-arrow map
  \[
    \ell_B:
    \Sigma\times L
    \longrightarrow
    \operatorname{dom}(B)_1
  \]
  satisfying
  \[
    t\,\ell_B
    =
    r_B\rho_B(1_\Sigma\times o_L),
  \]
  \[
    s\,\ell_B
    =
    r_B\rho_B(1_\Sigma\times d_L),
  \]
  and
  \[
    \lambda_B
    \bigl(
      \ell_B(\sigma,l),
      \rho_B(\sigma,o_Ll)
    \bigr)
    =
    \rho_B(\sigma,d_Ll).
  \]
\end{enumerate}
\end{definition}

The categorical arrow \(\ell_B(\sigma,l)\) runs from the anchor at the destination history to the
anchor at the origin history. Target-to-source processing therefore transports the state in the
operational direction from \(o_L(l)\) to \(d_L(l)\).

\begin{lemma}[Transport along legal history paths]
\label{lem:history-path-transport}
Let
\[
  l_1:h_0\to h_1,
  \ldots,
  l_n:h_{n-1}\to h_n
\]
be a composable legal history path. Then
\[
  \lambda_B
  \Bigl(
    \ell_B(l_1)\circ\cdots\circ\ell_B(l_n),
    \rho_B(\sigma,h_0)
  \Bigr)
  =
  \rho_B(\sigma,h_n).
\]
\end{lemma}

\begin{proof}
The result is immediate for the empty path. For the induction step, use the one-step synchronization
equation and the right-action multiplication law. The target-to-source ordering is exactly the
displayed categorical composite.
\end{proof}

For a generalized incumbent
\[
  \sigma:U\to\Sigma
\]
and history
\[
  h:U\to H,
\]
pullback gives synchronized selected states over \(U\).

\begin{definition}[Pointed and generated residual strategic games]
\label{def:residual-strategic-game}
The pointed residual strategic game at
\[
  (\sigma,h):U\to\Sigma\times H
\]
is the pullback of the strategic implementation to \(\mathcal E/U\), together with the three
selected ready states
\[
  \rho_F(\sigma,h),
  \qquad
  \rho_G(\sigma,h),
  \qquad
  \rho_E(\sigma,h).
\]

Let
\[
  i_{\sigma,h}:
  \mathbbm1_U+\mathbbm1_U+\mathbbm1_U
  \longrightarrow
  P_{\mathsf{SProt}}|_U
\]
be the induced map into the complete protocol carrier, using the branch carrier inclusions.

The generated residual strategic protocol is
\[
  \operatorname{Res}_{\sigma,h}
  \bigl(
    \mathsf{SProt}(\mathcal G)
  \bigr)
  :=
  \operatorname{Reach}
  \bigl(
    \mathsf{SProt}(\mathcal G)|_U,
    i_{\sigma,h}
  \bigr),
\]
equipped with total initialization.
\end{definition}

The branch carriers of the generated residual protocol are the pullbacks
\[
  P_B^{\mathrm{res}}
  :=
  P_B|_U
  \times_{P_{\mathsf{SProt}}|_U}
  \operatorname{Reach}_0.
\]
Since the original branch carriers are stable under their branch arrows, these intersections remain
stable branch restrictions.

The pointed residual implementation records the selected synchronized states. The generated residual
protocol is their transition-closed reachable realization. In particular, forward states may
generate backward states through return arrows inside the same complete protocol.

\subsection{Finite internal extensive games}

Fix an external natural number \(N\). A finite deterministic perfect-information game presentation
consists of:

\begin{enumerate}
  \item a finite constant object \(I\) of players;
  \item node objects \(H_0,\ldots,H_N\), with a root
  \[
    h_0:\mathbbm1\to H_0;
  \]
  \item for each \(n<N\), an object of legal move occurrences
  \[
    \pi_n:A_n\to H_n
  \]
  and a successor map
  \[
    \nu_n:A_n\to H_{n+1};
  \]
  \item an active-player map
  \[
    \iota_n:H_n\to I;
  \]
  \item a terminal outcome map
  \[
    u:H_N\to Z;
  \]
  \item for every player \(i\), an internal preference relation
  \[
    \succeq_i\hookrightarrow Z\times Z.
  \]
\end{enumerate}

The node objects may identify positions reached by different move sequences. The history compiler
therefore uses the canonical unfolding of the rooted game graph.

\subsection{The unfolded legal-history object}

Define
\[
  \widehat H_0:=\mathbbm1
\]
and
\[
  e_0:=h_0:\widehat H_0\to H_0.
\]

Recursively, define
\[
  \widehat H_{n+1}
  :=
  \widehat H_n
  \times_{e_n,H_n,\pi_n}
  A_n
\]
and
\[
  e_{n+1}
  :=
  \nu_n\pi_2:
  \widehat H_{n+1}\to H_{n+1}.
\]

The predecessor map is
\[
  \partial_n:=\pi_1:
  \widehat H_{n+1}\to\widehat H_n.
\]

Put
\[
  \widehat H
  :=
  \coprod_{n=0}^{N}\widehat H_n
\]
and
\[
  \widehat L
  :=
  \coprod_{n<N}\widehat H_{n+1}.
\]
On the \(n\)-th summand, the origin map is the inclusion of
\[
  \partial_n:\widehat H_{n+1}\to\widehat H_n
\]
and the destination map is the inclusion of \(\widehat H_{n+1}\) itself.

For \(n\leq m\), write
\[
  \partial_{n,m}:
  \widehat H_m\to\widehat H_n
\]
for the iterated predecessor map.

\begin{proposition}[Unique unfolded history words]
\label{prop:unique-unfolded-history-words}
For each \(n\), the object of length-\(n\) legal extension paths beginning at the root is canonically
isomorphic to \(\widehat H_n\). Consequently every element of \(\widehat H\) has a unique finite
legal extension word from the root.
\end{proposition}

\begin{proof}
For \(n=0\), both objects are terminal.

Suppose the statement holds at depth \(n\). A path of length \(n+1\) is a path of length \(n\)
together with a legal move occurrence at its endpoint. By the induction isomorphism, its object is
\[
  \widehat H_n
  \times_{H_n}
  A_n
  =
  \widehat H_{n+1}.
\]
The recursive construction records the predecessor at each stage, so the resulting extension word
is unique.
\end{proof}

The active-player and terminal maps pull back to
\[
  \widehat\iota_n:=\iota_ne_n,
  \qquad
  \widehat u:=ue_N.
\]

The legal move family over an unfolded node is
\[
  \widehat A_n
  :=
  \widehat H_n\times_{H_n}A_n
  =
  \widehat H_{n+1}
  \longrightarrow
  \widehat H_n.
\]

\subsection{Continuation strategy bundles}

For every player \(i\), let
\[
  \widehat H_{n,i}\hookrightarrow\widehat H_n
\]
be the subobject on which
\[
  \widehat\iota_n=i.
\]

For a fixed depth \(n\), define the object of future \(i\)-decision nodes over
\(\widehat H_n\) by
\[
  D_{i,n}
  :=
  \coprod_{m=n}^{N-1}\widehat H_{m,i},
\]
where the structure map to \(\widehat H_n\) on the \(m\)-th summand is
\[
  \partial_{n,m}.
\]

Let
\[
  B_{i,n}\to D_{i,n}
\]
be the family obtained by restricting
\[
  \widehat A_m\to\widehat H_m
\]
to \(\widehat H_{m,i}\) on each summand.

Define the continuation strategy object of player \(i\), as an object of
\(\mathcal E/\widehat H_n\), by
\[
  \Sigma_{i,n}
  :=
  \Pi_{D_{i,n}\to\widehat H_n}
  (B_{i,n}).
\]

Its fibre at an unfolded history \(\widehat h\) is the object of choices at every future
\(i\)-controlled decision node descending from \(\widehat h\).

Define the continuation-profile object
\[
  \Sigma_n
  :=
  \prod_{i\in I}^{\mathcal E/\widehat H_n}
  \Sigma_{i,n}.
\]

At the root,
\[
  \Sigma_i:=\Sigma_{i,0},
  \qquad
  \Sigma:=\prod_{i\in I}\Sigma_i.
\]

Restriction of a global strategy section to the descendant decision nodes gives, by
Beck--Chevalley and the universal property of dependent products, a morphism
\[
  \operatorname{res}_{i,n}:
  \widehat H_n\times\Sigma_i
  \longrightarrow
  \Sigma_{i,n}.
\]
Taking the product over players gives
\[
  \operatorname{res}_n:
  \widehat H_n\times\Sigma
  \longrightarrow
  \Sigma_n.
\]

\subsection{Continuation outcomes}

At depth \(N\), define
\[
  o_N:\Sigma_N\to Z
\]
over \(\widehat H_N\) by
\[
  o_N(\widehat h,*)=\widehat u(\widehat h).
\]

For \(n<N\), evaluation of the active player's continuation strategy at the current node gives a
legal move occurrence and hence a successor map
\[
  s_n:\Sigma_n\to\widehat H_{n+1}
\]
over \(\widehat H_n\).

Restriction of all continuation strategies to descendants of that successor gives a map
\[
  r_n:\Sigma_n\to\Sigma_{n+1}
\]
lying over \(s_n\).

Define recursively
\[
  o_n:=o_{n+1}r_n.
\]

Thus the fibre of
\[
  o_n:\Sigma_n\to Z
\]
at \(\widehat h\) is the continuation-outcome map
\[
  o_{\widehat h}:\Sigma_{\widehat h}\to Z
\]
for the continuation game rooted at \(\widehat h\).

For a continuation profile
\[
  \bar\sigma\in\Sigma_{\widehat h}
\]
and a candidate response profile
\[
  \bar\tau\in\Sigma_{\widehat h},
\]
define
\[
\begin{aligned}
  \chi_{\widehat h}(\bar\sigma,\bar\tau)
  :=
  \bigwedge_{i\in I}
  \forall_{\theta_i\in\Sigma_{i,\widehat h}}
  \Bigl[
    o_{\widehat h}(\bar\sigma_{-i},\bar\tau_i)
    \succeq_i
    o_{\widehat h}(\bar\sigma_{-i},\theta_i)
  \Bigr].
\end{aligned}
\]

The finite meet and the displayed universal quantifiers are internal operations in the topos.

Define the global history-indexed classifier
\[
\begin{aligned}
  \chi_{\mathrm{hist}}:
  \Sigma\times\widehat H\times\Sigma
  &\longrightarrow
  \Omega,
  \\
  \chi_{\mathrm{hist}}(\sigma,\widehat h,\tau)
  &:=
  \chi_{\widehat h}
  \bigl(
    \operatorname{res}_{\widehat h}\sigma,
    \operatorname{res}_{\widehat h}\tau
  \bigr).
\end{aligned}
\]

Its canonical extensional evaluator is the history-indexed response evaluator.

\begin{definition}[Deviation-complete unfolded-history compiler]
\label{def:deviation-complete-history-implementation}
The deviation-complete unfolded-history compiler of a finite deterministic perfect-information game
is the history-indexed implementation whose:
\begin{enumerate}
  \item history object is the full unfolded legal-history object \(\widehat H\);
  \item legal extension object is \(\widehat L\);
  \item synchronized ready states are those obtained by executing the compiled continuation arena
  and evaluator at each pair \((\sigma,\widehat h)\);
  \item branch extension maps are the compiled input arrows representing the corresponding legal
  unfolded extension;
  \item evaluator classifier is \(\chi_{\mathrm{hist}}\).
\end{enumerate}
\end{definition}

Every legal move occurrence is represented independently of the incumbent profile, and every
unfolded legal history is reachable from the root by its unique finite extension word.

\begin{definition}[Residual stability]
\label{def:residual-stability}
A strategy profile \(\sigma\) is residually stable when
\[
  \chi_{\widehat h}
  \bigl(
    \operatorname{res}_{\widehat h}\sigma,
    \operatorname{res}_{\widehat h}\sigma
  \bigr)
  =
  \top
\]
for every unfolded legal history \(\widehat h\).
\end{definition}

\begin{lemma}[Continuation equilibrium adequacy]
\label{lem:continuation-equilibrium-adequacy}
For every unfolded legal history \(\widehat h\), one has
\[
  \chi_{\widehat h}
  \bigl(
    \operatorname{res}_{\widehat h}\sigma,
    \operatorname{res}_{\widehat h}\sigma
  \bigr)
  =
  \top
\]
if and only if the restriction of \(\sigma\) to the continuation game rooted at \(\widehat h\) is a
Nash equilibrium of that continuation game.
\end{lemma}

\begin{proof}
Expanding the classifier on the diagonal gives
\[
\begin{aligned}
  \bigwedge_{i\in I}
  \forall_{\theta_i\in\Sigma_{i,\widehat h}}
  \Bigl[
    o_{\widehat h}(\operatorname{res}_{\widehat h}\sigma)
    \succeq_i
    o_{\widehat h}
    \bigl(
      (\operatorname{res}_{\widehat h}\sigma)_{-i},
      \theta_i
    \bigr)
  \Bigr].
\end{aligned}
\]
This states exactly that no player has a profitable unilateral continuation-strategy deviation in
the continuation game rooted at \(\widehat h\). That is the internal Nash equilibrium condition for
that continuation game.
\end{proof}

\begin{theorem}[Subgame perfection as residual stability]
\label{thm:subgame-perfection}
For a finite deterministic perfect-information game compiled by the deviation-complete
unfolded-history compiler, a strategy profile is subgame-perfect if and only if it is residually
stable.
\end{theorem}

\begin{proof}
Every legal move sequence, including every off-path sequence, is represented by a unique element of
\(\widehat H\). By Lemma~\ref{lem:continuation-equilibrium-adequacy}, residual stability at
\(\widehat h\) is equivalent to Nash equilibrium of the profile restriction in the continuation
game rooted at \(\widehat h\). Requiring this at every unfolded legal history is precisely
subgame-perfect equilibrium.
\end{proof}

\section{Phase-decorated execution spans}
\label{sec:phase-indexed-execution-spans}

Feedback is applied after relative execution evaluation. The exposed optic residual and the hidden
feedback object are different structures. Persistent machine state remains part of the visible
execution configuration. Only private event-routing objects are hidden.

Fix an object \(\Phi_0\) of phase labels.

\begin{definition}[Phase-decorated execution object]
\label{def:phase-indexed-execution-object}
A phase-decorated execution object is a pair
\[
  (X,\varphi_X)
\]
with
\[
  \varphi_X:X\to\Phi_0.
\]
\end{definition}

\begin{definition}[Phase-decorated execution span]
\label{def:phase-decorated-execution-span}
A phase-decorated execution span
\[
  (X,\varphi_X)\longrightarrow(Y,\varphi_Y)
\]
is an ordinary span
\[
  X\xleftarrow{\ell}U\xrightarrow{r}Y.
\]
An execution witness carries source phase \(\varphi_X\ell\) and target phase \(\varphi_Yr\).
Composition is ordinary pullback composition.
\end{definition}

The phase labels in this definition are endpoint decorations. They do not impose a separate
phase-transition legality relation.

Passing to isomorphism classes gives a coproduct-monoidal category
\[
  \operatorname{Span}^{\Phi,+}_{\omega,0}(\mathcal E).
\]
Because \(\mathcal E\) is a Grothendieck topos, it satisfies the feedback span-base conditions of
Chapter~9.

Let
\[
  U:X+H\longrightarrow Y+H
\]
be a phase-decorated span. By extensivity, it has blocks
\[
  U_{XY}:X\to Y,
  \qquad
  U_{XH}:X\to H,
\]
\[
  U_{HY}:H\to Y,
  \qquad
  U_{HH}:H\to H.
\]

\begin{definition}[Phase-decorated execution trace]
\label{def:phase-indexed-execution-trace}
The trace over \(H\) is
\[
  \operatorname{Tr}^{H}(U)
  :=
  U_{XY}
  \boxplus
  U_{HY}
  \circ
  \operatorname{Path}(U_{HH})
  \circ
  U_{XH}.
\]
\end{definition}

Phase labels remain on the visible boundary objects after the intermediate route object is hidden.

\section{Private linear subdivision}
\label{sec:private-linear-subdivision}

Let
\[
  M=\boxplus_{e\in E}M_e:X\longrightarrow Y
\]
be a finite homwise coproduct of macro-event execution spans.

For every \(e\), choose an integer \(n_e\geq0\). If \(n_e=0\), choose only a span
\[
  U_{e,0}:X\to Y
\]
with a specified isomorphism
\[
  U_{e,0}\cong M_e.
\]

If \(n_e\geq1\), choose private hidden objects
\[
  H_{e,1},\ldots,H_{e,n_e}
\]
and a chain
\[
  X
  \xrightarrow{U_{e,0}}
  H_{e,1}
  \xrightarrow{U_{e,1}}
  \cdots
  \xrightarrow{U_{e,n_e-1}}
  H_{e,n_e}
  \xrightarrow{U_{e,n_e}}
  Y
\]
with a specified span isomorphism
\[
  U_{e,n_e}\circ\cdots\circ U_{e,0}
  \cong
  M_e.
\]

Assume that hidden objects belonging to different macro events are disjoint and that no subdivision
arrow belongs to two chains.

\begin{definition}[Private linear subdivision]
\label{def:private-linear-subdivision}
Put
\[
  H
  :=
  \coprod_{e:n_e\geq1}
  \coprod_{i=1}^{n_e}
  H_{e,i}.
\]

The private linear subdivision of \(M\) is the span
\[
  \widetilde M:
  X+H
  \longrightarrow
  Y+H
\]
whose blocks are:
\[
  (\widetilde M)_{XY}
  :=
  \boxplus_{e:n_e=0}U_{e,0},
\]
\[
  (\widetilde M)_{XH}
  :=
  \boxplus_{e:n_e\geq1}U_{e,0},
\]
\[
  (\widetilde M)_{HH}
  :=
  \boxplus_{e:n_e\geq1}
  \boxplus_{i=1}^{n_e-1}
  U_{e,i},
\]
and
\[
  (\widetilde M)_{HY}
  :=
  \boxplus_{e:n_e\geq1}U_{e,n_e},
\]
with each span inserted into the block determined by its tagged source and target summands.
\end{definition}

A map of private linear subdivisions preserves the macro-event index, the ordered stage of every
hidden object, and every subdivision link.

\begin{theorem}[Finite subdivision collapse]
\label{thm:finite-subdivision-collapse}
For a private linear subdivision,
\[
  \operatorname{Tr}^{H}(\widetilde M)
  \cong
  M.
\]
The isomorphism is natural in maps of private linear subdivisions.
\end{theorem}

\begin{proof}
Write
\[
  R_{XY}:=(\widetilde M)_{XY},
  \qquad
  R_{XH}:=(\widetilde M)_{XH},
\]
\[
  R_{HH}:=(\widetilde M)_{HH},
  \qquad
  R_{HY}:=(\widetilde M)_{HY}.
\]

Because the hidden objects are tagged disjointly, a pullback between a span ending at \(H_{e,i}\)
and one beginning at \(H_{e',j}\) is initial unless
\[
  e=e',
  \qquad
  j=i+1.
\]
Consequently
\[
  R_{HY}\circ R_{HH}^{[m]}\circ R_{XH}
\]
is initial unless
\[
  m=n_e-1
\]
for one of the chains with \(n_e\geq1\). In that case the surviving summand is
\[
  U_{e,n_e}\circ\cdots\circ U_{e,0}.
\]

Therefore
\[
\begin{aligned}
  \operatorname{Tr}^{H}(\widetilde M)
  &=
  R_{XY}
  \boxplus
  \boxplus_{m\geq0}
  R_{HY}\circ R_{HH}^{[m]}\circ R_{XH}
  \\
  &\cong
  \boxplus_{e:n_e=0}U_{e,0}
  \boxplus
  \boxplus_{e:n_e\geq1}
  U_{e,n_e}\circ\cdots\circ U_{e,0}
  \\
  &\cong
  \boxplus_eM_e
  \\
  &=M.
\end{aligned}
\]

A map of subdivisions sends the unique surviving path of each source chain to the unique surviving
path of the corresponding target chain. Naturality follows from functoriality of span pullback and
path composition.
\end{proof}

Acyclicity alone is insufficient: an acyclic hidden graph may branch or merge and consequently
produce several hidden paths. Private linearity is the exact property used in the proof.

\section{Relative execution factorization of Mealy pipelines}
\label{sec:relative-execution-factorization}

Let
\[
  P:A\rightsquigarrow B,
  \qquad
  Q:B\rightsquigarrow C,
\]
and let
\[
  (Y\xrightarrow{r}C_0,\sigma_Y)
\]
be a right internal \(C\)-action.

The relative state object of the pipeline is
\[
\begin{aligned}
  S_{Q\circ P,Y}
  &:=
  \bigl(
    P_0\times_{q_P,B_0,p_Q}Q_0
  \bigr)
  \times_{q_Q,C_0,r}
  Y.
\end{aligned}
\]
Write an element as
\[
  (x,z,y)
\]
with
\[
  q_P(x)=p_Q(z),
  \qquad
  q_Q(z)=r(y).
\]

The primitive execution object is
\[
  M_{Q\circ P,Y}
  :=
  A_1
  \times_{t_A,A_0,p_P\pi_P}
  S_{Q\circ P,Y}.
\]
Write an element as
\[
  (a,x,z,y)
\]
with
\[
  t_A(a)=p_P(x).
\]

Define the intermediate handshake object
\[
\begin{aligned}
  H_{Q,P;Y}
  &:=
  P_0
  \times_{q_P,B_0,s_B}
  B_1
  \times_{t_B,B_0,p_Q}
  Q_0
  \times_{q_Q,C_0,r}
  Y.
\end{aligned}
\]
Its elements are written
\[
  (x',b,z,y)
\]
with
\[
  q_P(x')=s_B(b),
  \qquad
  t_B(b)=p_Q(z),
  \qquad
  q_Q(z)=r(y).
\]

Define
\[
\begin{aligned}
  h_0:
  M_{Q\circ P,Y}
  &\longrightarrow
  H_{Q,P;Y},
  \\
  h_0(a,x,z,y)
  &:=
  \bigl(
    \delta_P(a,x),
    \omega_P(a,x),
    z,
    y
  \bigr),
\end{aligned}
\]
and
\[
\begin{aligned}
  h_1:
  H_{Q,P;Y}
  &\longrightarrow
  S_{Q\circ P,Y},
  \\
  h_1(x',b,z,y)
  &:=
  \Bigl(
    x',
    \delta_Q(b,z),
    \sigma_Y(\omega_Q(b,z),y)
  \Bigr).
\end{aligned}
\]

The first map is well typed because
\[
  q_P\delta_P(a,x)
  =
  s_B\omega_P(a,x)
\]
and
\[
  t_B\omega_P(a,x)
  =
  q_P(x)
  =
  p_Q(z).
\]
The second is well typed because
\[
  p_Q\delta_Q(b,z)
  =
  s_B(b)
  =
  q_P(x')
\]
and
\[
  q_Q\delta_Q(b,z)
  =
  s_C\omega_Q(b,z)
  =
  r\,\sigma_Y(\omega_Q(b,z),y).
\]

Define spans
\[
  U_0
  :=
  \left(
    S_{Q\circ P,Y}
    \xleftarrow{s_{Q\circ P,Y}}
    M_{Q\circ P,Y}
    \xrightarrow{h_0}
    H_{Q,P;Y}
  \right)
\]
and
\[
  U_1
  :=
  \left(
    H_{Q,P;Y}
    \xleftarrow{1}
    H_{Q,P;Y}
    \xrightarrow{h_1}
    S_{Q\circ P,Y}
  \right).
\]

\begin{theorem}[Pipeline execution factorization]
\label{thm:pipeline-execution-factorization}
There is a canonical isomorphism
\[
  U_1\circ U_0
  \cong
  m_{Q\circ P,Y},
\]
where \(m_{Q\circ P,Y}\) is the Chapter~6 primitive execution span of the pipeline.

Let
\[
  \widetilde m_{Q,P;Y}:
  S_{Q\circ P,Y}+H_{Q,P;Y}
  \longrightarrow
  S_{Q\circ P,Y}+H_{Q,P;Y}
\]
have only the blocks
\[
  (\widetilde m_{Q,P;Y})_{SH}=U_0,
  \qquad
  (\widetilde m_{Q,P;Y})_{HS}=U_1.
\]
Then
\[
  \boxed{
  m_{Q\circ P,Y}
  \cong
  \operatorname{Tr}^{H_{Q,P;Y}}
  \bigl(
    \widetilde m_{Q,P;Y}
  \bigr).
  }
\]
\end{theorem}

\begin{proof}
Since the left leg of \(U_1\) is the identity, the pullback apex of
\[
  U_1\circ U_0
\]
is canonically \(M_{Q\circ P,Y}\). Its source leg is
\[
  (a,x,z,y)\longmapsto(x,z,y).
\]
Its target leg is
\[
\begin{aligned}
  h_1h_0(a,x,z,y)
  &=
  h_1
  \bigl(
    \delta_P(a,x),
    \omega_P(a,x),
    z,
    y
  \bigr)
  \\
  &=
  \Bigl(
    \delta_P(a,x),
    \delta_Q(\omega_P(a,x),z),
    \\
    &\hspace{28mm}
    \sigma_Y
    \bigl(
      \omega_Q(\omega_P(a,x),z),
      y
    \bigr)
  \Bigr).
\end{aligned}
\]

Pipeline composition in \(\operatorname{Mly}(\mathcal E)\) has
\[
  \delta_{Q\circ P}(a,(x,z))
  =
  \bigl(
    \delta_P(a,x),
    \delta_Q(\omega_P(a,x),z)
  \bigr)
\]
and
\[
  \omega_{Q\circ P}(a,(x,z))
  =
  \omega_Q(\omega_P(a,x),z).
\]
The target leg is therefore exactly the Chapter~6 primitive target map of the pipeline.

The \(HH\)-block and direct \(SS\)-block of \(\widetilde m_{Q,P;Y}\) are zero. Hence
\[
\begin{aligned}
  \operatorname{Tr}^{H_{Q,P;Y}}
  (\widetilde m_{Q,P;Y})
  &=
  U_1
  \circ
  \operatorname{Path}(0_{H_{Q,P;Y}})
  \circ
  U_0
  \\
  &\cong
  U_1\circ1_{H_{Q,P;Y}}\circ U_0
  \\
  &\cong
  m_{Q\circ P,Y}.
\end{aligned}
\]
\end{proof}

The visible object \(S_{Q\circ P,Y}\) retains the current \(P\)-state, current \(Q\)-state, and
downstream state. Only the intermediate emitted \(B\)-arrow and its handshake phase are hidden.

\section{Event realization of stateless maps}
\label{sec:event-realization}

A map-representable Mealy morphism between discrete categories has only identity input arrows.
Its Chapter~6 primitive execution span therefore does not directly recover the graph span of its
object map. A free event category supplies the missing input event.

\begin{definition}[Internal event category]
\label{def:internal-event-category}
Let
\[
  o,d:E\rightrightarrows W
\]
be an internal family of directed execution events, where \(o\) is the execution origin and \(d\) is
the execution destination.

Define
\[
  \mathsf{Ev}(o,d)
\]
to be the chosen free internal category on the graph whose object object is \(W\), whose generating
arrow object is \(E\), and whose categorical source and target are
\[
  s_{\mathsf{Ev}}=d,
  \qquad
  t_{\mathsf{Ev}}=o.
\]
\end{definition}

Thus a categorical generator is written
\[
  d(e)\longrightarrow o(e).
\]
Under target-to-source processing, it executes in the desired direction
\[
  o(e)\longrightarrow d(e).
\]

For an internal category \(A\), let
\[
  \mathsf{ObAct}_A
  :=
  (A_0\xrightarrow{1_{A_0}}A_0,s_A)
\]
be the canonical right action of \(A\) on its object object.

\begin{lemma}[Event realization]
\label{lem:event-realization}
Let
\[
  \iota:E\hookrightarrow\mathsf{Ev}(o,d)_1
\]
be the inclusion of length-one generators. Evaluate the identity Mealy morphism
\[
  1_{\mathsf{Ev}(o,d)}
\]
against the canonical object action
\[
  \mathsf{ObAct}_{\mathsf{Ev}(o,d)}.
\]
Restrict its primitive execution span along \(\iota\). The resulting span is
\[
  \boxed{
  W\xleftarrow{o}E\xrightarrow{d}W.
  }
\]
\end{lemma}

\begin{proof}
For the identity Mealy morphism on an internal category \(A\), evaluated at the canonical object
action, the relative state object is canonically
\[
  S_{1_A,\mathsf{ObAct}_A}\cong A_0,
\]
and the primitive execution object is canonically
\[
  M_{1_A,\mathsf{ObAct}_A}\cong A_1.
\]
The primitive source and target maps are
\[
  t_A,
  \qquad
  s_A,
\]
respectively. Restriction along \(\iota\) therefore gives
\[
  t_A\iota=o,
  \qquad
  s_A\iota=d.
\]
\end{proof}

\begin{corollary}[Graph spans as primitive event execution]
\label{cor:graph-spans-as-event-execution}
Let
\[
  f:X\to Y.
\]
Put
\[
  W:=X+Y,
  \qquad
  E:=X,
\]
and define
\[
  o:=\iota_X,
  \qquad
  d:=\iota_Yf.
\]
Then the \(X\)-to-\(Y\) block of the restricted primitive execution span of
\(\mathsf{Ev}(o,d)\) is
\[
  \Gamma(f):
  X\xleftarrow{1_X}X\xrightarrow{f}Y.
\]
\end{corollary}

\begin{proof}
Apply Lemma~\ref{lem:event-realization} and take the block from the \(X\)-summand of \(W\) to the
\(Y\)-summand.
\end{proof}

\section{Open-game arena composition as step execution and feedback}
\label{sec:open-game-feedback}

For an open-normal game
\[
  \mathcal K=(\Lambda,P_{\mathcal K},C_{\mathcal K},\chi_{\mathcal K})
  :
  (X,S)\to(Z,Q),
\]
define
\[
  V_F^{\mathrm{in}}(\mathcal K)
  :=
  \Lambda\times X,
\]
\[
  V_F^{\mathrm{out}}(\mathcal K)
  :=
  \Lambda\times X\times Z,
\]
and
\[
\begin{aligned}
  \bar P_{\mathcal K}:
  V_F^{\mathrm{in}}(\mathcal K)
  &\longrightarrow
  V_F^{\mathrm{out}}(\mathcal K),
  \\
  (\lambda,x)
  &\longmapsto
  (\lambda,x,P_{\mathcal K}(\lambda,x)).
\end{aligned}
\]

Define
\[
  V_G^{\mathrm{in}}(\mathcal K)
  :=
  \Lambda\times X\times Q,
\]
\[
  V_G^{\mathrm{out}}(\mathcal K)
  :=
  \Lambda\times X\times Q\times S,
\]
and
\[
\begin{aligned}
  \bar C_{\mathcal K}:
  V_G^{\mathrm{in}}(\mathcal K)
  &\longrightarrow
  V_G^{\mathrm{out}}(\mathcal K),
  \\
  (\lambda,x,q)
  &\longmapsto
  (\lambda,x,q,C_{\mathcal K}(\lambda,x,q)).
\end{aligned}
\]

\begin{definition}[Extracted arena protocol and step execution]
\label{def:step-execution-open-normal}
Define
\[
  \mathsf{Prot}_{\mathrm{ar}}(\mathcal K)
  :=
  (\bar P_{\mathcal K},\bar C_{\mathcal K}).
\]
Its step execution span is
\[
  \operatorname{Exec}_{\mathrm{step}}
  \bigl(
    \mathsf{Prot}_{\mathrm{ar}}(\mathcal K)
  \bigr)
  :=
  \Gamma(\bar P_{\mathcal K})
  +
  \Gamma(\bar C_{\mathcal K}).
\]
\end{definition}

This definition is independent of the feedback construction below.

Now let
\[
  \mathcal G:(X,S)\to(Y,R),
  \qquad
  \mathcal H:(Y,R)\to(Z,Q)
\]
be open-normal games with strategy objects \(\Sigma,T\).

Take the phase object
\[
  \Phi_0
  :=
  \mathbbm1
  +
  \mathbbm1
  +
  \mathbbm1
  +
  \mathbbm1
  +
  \mathbbm1
  +
  \mathbbm1
\]
with its six coproduct injections denoted
\[
  \mathsf f_{\mathrm{in}},
  \mathsf f_{\mathrm{mid}},
  \mathsf f_{\mathrm{out}},
  \mathsf g_{\mathrm{in}},
  \mathsf g_{\mathrm{mid}},
  \mathsf g_{\mathrm{out}}
  :
  \mathbbm1\to\Phi_0.
\]

The objects below carry the corresponding constant phase maps.

\subsection{Forward subdivision}

Define
\[
  V_F^{\mathrm{in}}
  :=
  \Sigma\times T\times X,
\]
\[
  H_Y
  :=
  \Sigma\times T\times X\times Y,
\]
\[
  V_F^{\mathrm{out}}
  :=
  \Sigma\times T\times X\times Z.
\]
Define
\[
\begin{aligned}
  f_0:
  V_F^{\mathrm{in}}
  &\longrightarrow
  H_Y,
  \\
  (\sigma,\tau,x)
  &\longmapsto
  \bigl(
    \sigma,\tau,x,
    P_{\mathcal G}(\sigma,x)
  \bigr),
\end{aligned}
\]
and
\[
\begin{aligned}
  f_1:
  H_Y
  &\longrightarrow
  V_F^{\mathrm{out}},
  \\
  (\sigma,\tau,x,y)
  &\longmapsto
  \bigl(
    \sigma,\tau,x,
    P_{\mathcal H}(\tau,y)
  \bigr).
\end{aligned}
\]

\subsection{Backward subdivision}

Define
\[
  V_G^{\mathrm{in}}
  :=
  \Sigma\times T\times X\times Q,
\]
\[
  H_R
  :=
  \Sigma\times T\times X\times Y\times Q\times R,
\]
\[
  V_G^{\mathrm{out}}
  :=
  \Sigma\times T\times X\times Q\times S.
\]
Define
\[
\begin{aligned}
  g_0:
  V_G^{\mathrm{in}}
  &\longrightarrow
  H_R,
  \\
  (\sigma,\tau,x,q)
  &\longmapsto
  \Bigl(
    \sigma,
    \tau,
    x,
    P_{\mathcal G}(\sigma,x),
    q,
    C_{\mathcal H}
    \bigl(
      \tau,
      P_{\mathcal G}(\sigma,x),
      q
    \bigr)
  \Bigr),
\end{aligned}
\]
and
\[
\begin{aligned}
  g_1:
  H_R
  &\longrightarrow
  V_G^{\mathrm{out}},
  \\
  (\sigma,\tau,x,y,q,r)
  &\longmapsto
  \bigl(
    \sigma,
    \tau,
    x,
    q,
    C_{\mathcal G}(\sigma,x,r)
  \bigr).
\end{aligned}
\]

The first backward link computes the actual intermediate positive output
\[
  y=P_{\mathcal G}(\sigma,x)
\]
before applying the downstream coplay. Consequently the subdivision does not contain an independent
unconstrained \(y\)-coordinate.

\subsection{The route event category}

Put
\[
  V^{\mathrm{in}}
  :=
  V_F^{\mathrm{in}}
  +
  V_G^{\mathrm{in}},
\]
\[
  V^{\mathrm{out}}
  :=
  V_F^{\mathrm{out}}
  +
  V_G^{\mathrm{out}},
\]
\[
  H:=H_Y+H_R,
\]
and
\[
  W:=V^{\mathrm{in}}+H+V^{\mathrm{out}}.
\]

Define the micro-event object
\[
  E_{\mathrm{micro}}
  :=
  V_F^{\mathrm{in}}
  +
  H_Y
  +
  V_G^{\mathrm{in}}
  +
  H_R.
\]

Its origin and destination maps
\[
  o,d:E_{\mathrm{micro}}\rightrightarrows W
\]
are defined summandwise by
\[
\begin{array}{c|c|c}
\text{event summand}&o&d\\
\hline
V_F^{\mathrm{in}}
&
V_F^{\mathrm{in}}\hookrightarrow W
&
V_F^{\mathrm{in}}\xrightarrow{f_0}H_Y\hookrightarrow W
\\[1mm]
H_Y
&
H_Y\hookrightarrow W
&
H_Y\xrightarrow{f_1}V_F^{\mathrm{out}}\hookrightarrow W
\\[1mm]
V_G^{\mathrm{in}}
&
V_G^{\mathrm{in}}\hookrightarrow W
&
V_G^{\mathrm{in}}\xrightarrow{g_0}H_R\hookrightarrow W
\\[1mm]
H_R
&
H_R\hookrightarrow W
&
H_R\xrightarrow{g_1}V_G^{\mathrm{out}}\hookrightarrow W.
\end{array}
\]

Define the route event category
\[
  \mathsf{Route}_{\mathcal H,\mathcal G}
  :=
  \mathsf{Ev}(o,d).
\]

By Lemma~\ref{lem:event-realization}, the primitive execution span of the identity Mealy morphism on
\(\mathsf{Route}_{\mathcal H,\mathcal G}\), restricted to the length-one generators, is the
phase-decorated span
\[
  \widetilde U_{\mathcal H,\mathcal G}:
  V^{\mathrm{in}}+H
  \longrightarrow
  V^{\mathrm{out}}+H
\]
whose nonzero blocks are
\[
  (\widetilde U_{\mathcal H,\mathcal G})_{V^{\mathrm{in}},H}
  =
  \Gamma(f_0)+\Gamma(g_0),
\]
and
\[
  (\widetilde U_{\mathcal H,\mathcal G})_{H,V^{\mathrm{out}}}
  =
  \Gamma(f_1)+\Gamma(g_1).
\]
Its direct visible block and hidden--hidden block are zero.

There are no arrows between \(H_Y\) and \(H_R\), and no additional hidden arrows.

\subsection{Composite branch equations}

\begin{lemma}[Extracted composite branches]
\label{lem:extracted-composite-branches}
The extracted forward and backward maps of the open-normal composite satisfy
\[
  \bar P_{\mathcal H\circ\mathcal G}
  =
  f_1f_0
\]
and
\[
  \bar C_{\mathcal H\circ\mathcal G}
  =
  g_1g_0.
\]
\end{lemma}

\begin{proof}
By definition,
\[
\begin{aligned}
  \bar P_{\mathcal H\circ\mathcal G}
  (\sigma,\tau,x)
  &=
  \Bigl(
    \sigma,\tau,x,
    P_{\mathcal H}
    \bigl(
      \tau,
      P_{\mathcal G}(\sigma,x)
    \bigr)
  \Bigr)
  \\
  &=
  f_1f_0(\sigma,\tau,x).
\end{aligned}
\]

Similarly,
\[
\begin{aligned}
  \bar C_{\mathcal H\circ\mathcal G}
  (\sigma,\tau,x,q)
  &=
  \Bigl(
    \sigma,\tau,x,q,
    C_{\mathcal G}
    \bigl(
      \sigma,
      x,
      C_{\mathcal H}
      (
        \tau,
        P_{\mathcal G}(\sigma,x),
        q
      )
    \bigr)
  \Bigr)
  \\
  &=
  g_1g_0(\sigma,\tau,x,q).
\end{aligned}
\]
\end{proof}

Inside the length-two path object
\[
  \mathsf{Route}_{\mathcal H,\mathcal G,1}^{[2]}
\]
there is a subobject
\[
  E_{\mathrm{macro}}
  :=
  V_F^{\mathrm{in}}
  +
  V_G^{\mathrm{in}}
\]
whose two families are the categorical composites whose operational orders are respectively
\[
  f_0\text{ then }f_1,
  \qquad
  g_0\text{ then }g_1.
\]
Restricting primitive execution to these macro-event arrows gives
\[
  \Gamma(f_1f_0)
  +
  \Gamma(g_1g_0).
\]

\begin{theorem}[Strategic arena composition as step execution feedback]
\label{thm:strategic-arena-composition-feedback}
There is a canonical isomorphism of phase-decorated execution spans
\[
  \boxed{
  \operatorname{Exec}_{\mathrm{step}}
  \bigl(
    \mathsf{Prot}_{\mathrm{ar}}
    (\mathcal H\circ\mathcal G)
  \bigr)
  \cong
  \operatorname{Tr}^{H}
  \bigl(
    \widetilde U_{\mathcal H,\mathcal G}
  \bigr).
  }
\]
The forward component is the graph span of
\[
  (\sigma,\tau,x)
  \longmapsto
  \bigl(
    \sigma,\tau,x,
    P_{\mathcal H}
    (
      \tau,
      P_{\mathcal G}(\sigma,x)
    )
  \bigr),
\]
and the backward component is the graph span of
\[
\begin{aligned}
  (\sigma,\tau,x,q)
  \longmapsto
  \Bigl(
    \sigma,\tau,x,q,
    C_{\mathcal G}
    \bigl(
      \sigma,
      x,
      C_{\mathcal H}
      (
        \tau,
        P_{\mathcal G}(\sigma,x),
        q
      )
    \bigr)
  \Bigr).
\end{aligned}
\]
\end{theorem}

\begin{proof}
The forward and backward route graphs are disjoint private linear chains. Their hidden--hidden block
is zero, so
\[
  \operatorname{Path}(0_H)\cong1_H.
\]
Hence
\[
\begin{aligned}
  \operatorname{Tr}^{H}
  (\widetilde U_{\mathcal H,\mathcal G})
  &\cong
  \bigl(
    \Gamma(f_1)+\Gamma(g_1)
  \bigr)
  \circ
  \bigl(
    \Gamma(f_0)+\Gamma(g_0)
  \bigr).
\end{aligned}
\]

Disjointness of \(H_Y\) and \(H_R\) makes the two cross-composites initial. Composition of graph
spans gives
\[
  \Gamma(f_1)\circ\Gamma(f_0)
  \cong
  \Gamma(f_1f_0)
\]
and
\[
  \Gamma(g_1)\circ\Gamma(g_0)
  \cong
  \Gamma(g_1g_0).
\]
Therefore
\[
  \operatorname{Tr}^{H}
  (\widetilde U_{\mathcal H,\mathcal G})
  \cong
  \Gamma(f_1f_0)+\Gamma(g_1g_0).
\]

By Lemma~\ref{lem:extracted-composite-branches},
\[
  \Gamma(f_1f_0)+\Gamma(g_1g_0)
  =
  \Gamma(\bar P_{\mathcal H\circ\mathcal G})
  +
  \Gamma(\bar C_{\mathcal H\circ\mathcal G}),
\]
which is
\[
  \operatorname{Exec}_{\mathrm{step}}
  \bigl(
    \mathsf{Prot}_{\mathrm{ar}}
    (\mathcal H\circ\mathcal G)
  \bigr)
\]
by Definition~\ref{def:step-execution-open-normal}.
\end{proof}

The response branch is not hidden by this trace. It is composed separately by stationary context
decomposition and meet. Adjoining that evaluator branch to the traced arena branch gives the
complete canonical protocol of the composite strategic game.

\begin{remark}[Persistent state in general strategic arenas]
For open-normal games, the intermediate value
\[
  P_{\mathcal G}(\sigma,x)
\]
may be recomputed from the visible strategy and input coordinates, so it may be placed in a private
hidden route object during the macro event.

For a genuinely stateful arena, the current Mealy carrier state cannot in general be reconstructed.
Theorem~\ref{thm:pipeline-execution-factorization} keeps the complete relative state object visible
and traces only the intermediate emitted-arrow handshake.
\end{remark}

\begin{remark}[No bicategorical \(\operatorname{Int}\) claim]
The preceding theorems are theorems about execution spans after relative evaluation and event
elaboration. They do not construct a traced symmetric monoidal bicategory of strategic Mealy systems,
nor a bicategorical \(\operatorname{Int}\) construction.

An ordinary \(\operatorname{Int}\)-morphism separates positive and negative span blocks. Additional
indexed visible-configuration structure would be required to prove that the backward block retains
the same persistent state generated by the forward block. The finite-path feedback theorems above do
not require that additional claim.
\end{remark}

\section{Examples}
\label{sec:strategic-mealy-examples}

\subsection{An internal selection-function decision}
\label{subsec:selection-function-decision}

Let \(Y,R\in\mathcal E\). An internal multivalued selection function may be represented by a map
\[
  \varepsilon:
  R^Y\longrightarrow\Omega^Y.
\]
Equivalently, it is a subobject
\[
  \mathsf{Sel}
  \hookrightarrow
  R^Y\times Y.
\]

Define
\[
  \mathcal D_\varepsilon:
  (\mathbbm1,\mathbbm1)
  \longrightarrow
  (Y,R)
\]
with strategy object
\[
  \Sigma=Y,
\]
play map
\[
  P(y,*)=y,
\]
unique coplay, and classifier
\[
  \chi_{\mathcal D_\varepsilon}
  (y,*,k,y')
  :=
  \operatorname{ev}
  \bigl(
    \varepsilon(k),
    y'
  \bigr).
\]
The equilibrium subobject is therefore classified by
\[
  (y,k)
  \longmapsto
  \operatorname{ev}
  \bigl(
    \varepsilon(k),
    y
  \bigr).
\]

The canonical residual is \(Y\), with forward map
\[
  y\longmapsto(y,y).
\]
The response evaluator may be the canonical extensional evaluator or a stateful controller that
constructs the complete outcome profile and then applies \(\varepsilon\).

\subsection{A sequential two-player game}
\label{subsec:sequential-two-player-game}

Let \(A\in\mathcal E\) be the first player's action object, and let
\[
  \pi_B:B\to A
\]
be an internally indexed family of second-player response objects. The contingent strategy object
of the second player is the dependent product
\[
  \Sigma_2
  :=
  \Pi_{!_A}(B).
\]
Its evaluation map is
\[
  \operatorname{ev}:
  \Sigma_2\times A\to B
\]
over \(A\).

Let
\[
  U:B\to\mathbb R_{\mathcal E}\times\mathbb R_{\mathcal E}
\]
be an internal payoff map, where \(\mathbb R_{\mathcal E}\) is a chosen internal ordered-real
object. Composite play is
\[
  (a,\tau)\longmapsto\operatorname{ev}(\tau,a).
\]

The first player's continuation against incumbent \(\tau\) is obtained by evaluating the payoff at
\[
  \operatorname{ev}(\tau,a')
\]
for every candidate \(a'\). The second player's local context is the unfolded history determined by
the first incumbent action. Thus the stationary context decomposition reconstructs the usual
normal-form response conditions.

The deviation-complete unfolded-history compiler includes every first-player action as a legal
history, whether or not it is selected by the incumbent first-player strategy. Residual stability
therefore requires the second-player strategy to be optimal in every fibre \(B_a\), including fibres
over off-path first-player actions. Theorem~\ref{thm:subgame-perfection} identifies this condition
with subgame perfection.

\subsection{Take-away and observation-relative minimization}
\label{subsec:take-away-minimization}

Let \(N\) and \(k\) be external positive integers and let the corresponding finite constant objects
be used internally. A nonterminal position is
\[
  (n,i),
\]
where \(n\) is a positive heap size and \(i\) is the active player. A legal move removes
\[
  m\in\{1,\ldots,k\},
  \qquad
  m\leq n,
\]
and sends
\[
  (n,i)\longmapsto(n-m,1-i).
\]

A full-history controller stores the complete path. A positional controller stores only the current
position. A residue controller stores
\[
  n\bmod(k+1).
\]

If the output protocol observes only eventual win or loss under a fixed residue policy, the
strategic behaviour specification may factor through the residue object, with the terminally
truncated heaps represented separately. The unrestricted arena protocol does not admit this
quotient: exact heap sizes have different enabled moves and different remaining path lengths.

Thus the residue quotient is a Nerode quotient of the chosen winner specification, not of the
unrestricted transition protocol. The distinction is forced by
Theorem~\ref{thm:minimal-strategic-realization}: minimization is always relative to the selected
ready-state behaviour map.

\subsection{Finite repeated prisoner's dilemma}
\label{subsec:repeated-prisoners-dilemma}

Let
\[
  A=\{C,D\}
\]
be the constant two-action object, and let the internal payoff map satisfy
\[
  T>R>P>S
\]
in the chosen ordered payoff object. Fix an external finite horizon \(N\).

The repeated arena stores:
\begin{enumerate}
  \item the current round;
  \item both controller states;
  \item the accumulated payoff.
\end{enumerate}
Each stage executes the synchronous protocol:
\begin{enumerate}
  \item both controllers emit from their current states;
  \item the joint action is evaluated;
  \item each controller observes the opponent's emitted action;
  \item both controller states update;
  \item the round counter advances.
\end{enumerate}

Tit-for-tat has state object \(\{C,D\}\), initialized at \(C\), emits its state, and updates to the
observed opponent action. Grim trigger has states
\[
  \{\mathsf{clean},\mathsf{triggered}\}
\]
with absorbing triggered state.

This system lies outside the open-normal fragment because deployed strategies are persistent
stateful controllers. It nevertheless lies in the general strategic Mealy bicategory. Product tensor
produces simultaneous stage output, relative execution supplies the common visible state object, and
Theorem~\ref{thm:pipeline-execution-factorization} factors each pipeline stage through a private
handshake object.

For a finite horizon, iterating those factorizations and applying the substitution and vanishing
equations of the Chapter~9 trace gives one trace over the coproduct of the finitely many private
stage-handshake objects. Controller states, round number, and accumulated payoff remain visible
throughout. Strategic residual semantics identifies histories precisely when they induce equal
future forward, backward, and counterfactual behaviours.

\section{Chapter summary}
\label{sec:strategic-mealy-games-summary}

This chapter constructed strategic interaction internally in a Grothendieck topos.

First, initialization was strengthened to a transition-closed right action of the input category.
This strengthening is forced by horizontal composition of output-lax Mealy simulations: the
downstream initialization state must be transported along the output-comparison arrow of the
upstream simulation. Transition-closed initialized Mealy morphisms form the symmetric monoidal
bicategory
\[
  \operatorname{IMly}^{\mathrm{cl}}(\mathcal E).
\]

Second, operational Mealy optics were constructed from initialized forward and backward legs.
Residual reparametrizers carry equivariant discard sections. These sections provide the comparison
\[
  (!_M)_\#R
  \Longrightarrow
  (!_N)_\#
\]
needed to make visible output functorial. Operational optics form
\[
  \operatorname{MOpt}_{\mathrm{op}}(\mathcal E).
\]

Third, pure map-representable residual slides were formulated as map-representable residual changes
equipped with strict invertible initialized leg comparisons and were identified by a homwise
coidentifier. This produced
\[
  \overline{\operatorname{MOpt}}_{\mathrm{op}}(\mathcal E).
\]
The quotient makes pure residual-presentation changes identities while retaining stateful residual
reparametrizations. The wide subcategories generated by raw tight cells form a locally wide
symmetric monoidal fragment.

Fourth, frozen strategy parameters were represented by discrete internal categories. Parametrized
arenas and contravariant parameter maps form
\[
  \overline{\operatorname{ParaMOpt}}_{\mathrm{op}}(\mathcal E).
\]
Generalized incumbents are handled by arbitrary base change rather than by global points. The
internally defined constructions are stable under base change through canonical pseudonatural
comparison isomorphisms.

Fifth, strategic context doctrines were formulated internally with explicitly typed
pseudonaturality, associativity, unit, symmetry, and middle-four comparison cells. Internal deviation
categories were constructed from free internal monoids, and internal query categories were formed
using the core groupoid of the context category.

Truth output is the one-object internal category associated to
\[
  (\Omega,\wedge,\top).
\]
Context-transparent evaluator machines transport along internal context isomorphisms. Response-strict
evaluator simulations preserve one-query classifiers exactly and are closed under the bicategorical
operations used in strategic composition.

These structures form the symmetric monoidal bicategory
\[
  \operatorname{StrMlyGame}^{=}_{\mathfrak C}(\mathcal K).
\]

Sixth, internal deterministic lenses were identified with tight discrete stateless
residual-codescent optics. This supplied the parametrized deterministic lens bicategory
\[
  \overline{\mathcal K}_{\mathrm{td}}.
\]
The stationary context doctrine was then constructed from the internal context object
\[
  X\times R^Y.
\]
Its sequential and parallel decomposition maps were derived by exponential transpose, and their
coherence was proved by equality of the resulting local continuation maps.

Seventh, internal deterministic open games were defined by strategy objects, play maps, coplay maps,
and best-response subobjects. The first bridge theorem proved
\[
  \boxed{
  \mathbf{OpenGame}_2(\mathcal E)
  \cong
  \mathbf{OGMly}^{\mathrm{nf}}_2(\mathcal E).
  }
\]
Open-normal games realize locally faithfully in the strategic Mealy bicategory based on the
stationary doctrine.

Eighth, structured canonical branch-faithful protocol machines were constructed using chosen free
internal protocol categories. Their forward, backward, and evaluator carrier summands are stable
branch restrictions. A strategy-port branch and tether generators force every structured protocol
simulation to use one common strategy reparametrization on all three branches.

For a strategy map \(u:T\to\Sigma\), the input reindexing functor always exists. The output
reindexing functor exists precisely when the play and coplay equations hold. After that functor has
been constructed, the carrier map into the input-pullback carrier defines a strict protocol
simulation precisely when the classifier equation holds.

Protocol composition was defined by extraction, strategic composition, and recompilation. The
second bridge theorem proved
\[
  \boxed{
  \mathbf{OGMly}^{\mathrm{nf}}_2(\mathcal E)
  \cong
  \mathbf{CanProt}^{\mathrm{str}}_2(\mathcal E).
  }
\]
Consequently,
\[
  \boxed{
  \mathbf{OpenGame}_2(\mathcal E)
  \cong
  \mathbf{CanProt}^{\mathrm{str}}_2(\mathcal E).
  }
\]

Forgetting only the literal chosen presentation while retaining branch, port, generator, and
ready-family structure yields the symmetric monoidal biequivalence
\[
  \mathbf{CanProt}^{\mathrm{str}}_2(\mathcal E)
  \simeq
  \mathbf{CanProt}^{\mathrm{rep}}_2(\mathcal E).
\]
No claim is made after forgetting all structural markings, because extraction and compiler
composition are then no longer determined by the underlying Mealy machine alone.

Ninth, the complete branch-faithful protocol was subjected to the residual semantics of Chapter~3.
Since a Grothendieck topos satisfies the standing exactness conditions, the strategic Nerode quotient
exists effectively and is canonically the derivative image. The resulting restricted canonical
behaviour machine is the unique reachable behaviourally separated minimal strategic realization.

On evaluator derivatives, complete finite-word response behaviour refines one-query response
observation. For canonical extensional evaluators, the two evaluator equivalence relations coincide.
No such collapse is asserted for the complete forward--backward--evaluator protocol.

Tenth, history-indexed synchronization was formulated over \(\Sigma\times H\) using explicitly typed
branch input-arrow maps for legal one-step extensions. A selected incumbent-history tuple defines a
pointed residual implementation; its generated residual protocol is the reachable transition-closed
submachine produced by the orbit-image construction.

For finite extensive games, the rooted game graph was canonically unfolded into path-history objects
\[
  \widehat H_n.
\]
Every unfolded history has a unique finite extension word from the root. Dependent continuation
strategy bundles, restriction maps, and continuation outcome maps were then constructed over the
unfolded history object. Residual stability was proved equivalent to subgame-perfect equilibrium.

Finally, feedback was performed after genuine relative execution evaluation.

For stateful Mealy pipelines, the primitive execution span factors through an intermediate emitted
arrow and satisfies
\[
  \boxed{
  m_{Q\circ P,Y}
  \cong
  \operatorname{Tr}^{H_{Q,P;Y}}
  \bigl(
    \widetilde m_{Q,P;Y}
  \bigr).
  }
\]
Persistent implementation and downstream state remain visible.

For stateless open-normal branches, a free internal event category turns an object map into a genuine
input event. Restriction of Chapter~6 primitive execution to the event generators recovers the graph
span of the map. Private linear subdivisions, including direct zero-hidden-stage events, collapse
under finite-path trace.

Step execution of an extracted arena protocol was defined independently as the coproduct tensor of
the graph spans of its forward and backward visible configuration maps. Sequential open-normal arena
composition then satisfies
\[
  \boxed{
  \operatorname{Exec}_{\mathrm{step}}
  \bigl(
    \mathsf{Prot}_{\mathrm{ar}}
    (\mathcal H\circ\mathcal G)
  \bigr)
  \cong
  \operatorname{Tr}^{H}
  \bigl(
    \widetilde U_{\mathcal H,\mathcal G}
  \bigr).
  }
\]
The response branch remains a separate strategic component composed by context decomposition and
meet.